\BourbakiExercisesSection

¶ 1. 设 A 为交换环，E、F 为两个 A-代数，M 为左 E-模，N 为左 F-模，G 为代数 \(\mathbf{G} = \mathbf{E} \otimes_{\mathbf{A}} \mathbf{F}\) 。对每一对有序的自同态 \(u \in \operatorname{End}_{\mathbf{E}}(\mathbf{M})\) , \(v \in \operatorname{End}_{\mathbf{F}}(\mathbf{N})\) ，存在唯一的 G-自同态 \(w\) （定义在 \(\mathbf{M} \otimes_{\mathbf{A}} \mathbf{N}\) 上），使得 \(w(x \otimes y) = u(x) \otimes v(y)\) 对所有 \(x \in \mathbf{M}\) , \(y \in \mathbf{N}\) 成立，并且我们可写 \(w = \phi(u \otimes v)\)，其中 \(\phi\) 是从 \(\operatorname{End}_{\mathbf{E}}(\mathbf{M}) \otimes_{\mathbf{A}} \operatorname{End}_{\mathbf{F}}(\mathbf{N})\) 到 \(\operatorname{End}_{\mathbf{G}}(\mathbf{M} \otimes_{\mathbf{A}} \mathbf{N})\) 中的 A-线性映射（称为典范映射）。以下设 A 为一个域。

\begin{enumerate}
\def\labelenumi{(\alph{enumi})}
\item
  证明典范映射 \(\phi\) 是单射（考虑元素 \(z = \sum_{i} u_{i} \otimes v_{i}\) ，其中 \(v_{i} \in \operatorname{End}_{\mathbf{F}}(\mathbf{N})\) 在 A 上线性无关，并写出 \((\phi(z))(x_{\alpha} \otimes y) = 0\) ，其中 \(x_{\alpha}\) 取遍 M 在 A 上的一个基的元素，而 \(y \in \mathbf{N}\) 任意）。
\item
  假设 \(\mathbf{M}\) 是有限生成的自由 E-模。证明在下列每种情形映射 \(\phi\) 都是双射：（1）E 在域 A 上是有限秩的；（2）N 是有限生成的 F-模。
\item
  假设 \(\mathbf{M}\) 是自由 E-模，且存在 \(y \in \mathbf{N}\) 与自同态的无限序列 \((v_n)\) （诸自同态定义在 F-模 \(\mathbf{N}\) 上），使得 \(\mathbf{N}\) 的由诸 \(v_n(y)\) 生成的（A 上的）向量子空间在 A 上具有无限秩。证明：若映射 \(\phi\) 是双射，则 \(\mathbf{M}\) 在 \(\mathbf{E}\) 上的每个基都必然是有限的。
\item
  仍设 \(\mathbf{M}\) 是自由 E-模。证明：G-模 \(\mathbf{M} \otimes_{\mathbf{A}} \mathbf{N}\) 是忠实的，当且仅当 F-模 \(\mathbf{N}\) 是忠实的（考虑 \(\mathbf{E}\) 在 \(\mathbf{A}\) 上的一个基）.
\end{enumerate}

\begin{enumerate}
\def\labelenumi{\arabic{enumi}.}
\setcounter{enumi}{1}
\tightlist
\item
  设 K 是交换域，E、F 是两个中心包含
\end{enumerate}

K 的域，M 是 E 上的左向量空间，N 是 F 上的左向量空间；设 \(G = E \otimes_{K} F\) 。

\begin{enumerate}
\def\labelenumi{(\alph{enumi})}
\item
  设 \(x \neq 0\) 是 \(\mathbf{M}\) 的一个元素，\(y \neq 0\) 是 \(\mathbf{N}\) 的一个元素。证明 \(x \otimes y\) 在 G-模 \(\mathbf{M} \otimes_{\mathbb{K}} \mathbf{N}\) 中是自由的（利用习题 1(d) 以及 \(\operatorname{End}_{\mathbf{E}}(\mathbf{M})\)（相应地 \(\operatorname{End}_{\mathbf{F}}(\mathbf{N})\)）在 \(\mathbf{M}\)（相应地 \(\mathbf{N}\)）中异于 0 的元素所成集合上的传递性）.
\item
  证明 \(\mathbf{M} \otimes_{\mathbb{K}} \mathbf{N}\) 是自由 G-模。（若 \((m_{\alpha})\) 是 \(\mathbf{M}\) 在 \(\mathbf{E}\) 上的一个基，\((n_{\beta})\) 是 \(\mathbf{N}\) 在 \(\mathbf{F}\) 上的一个基，则用与 (a) 类似的方法证明诸 G-模 \(\mathrm{G}(m_{\alpha} \otimes n_{\beta})\) 之和是直和。）
\end{enumerate}

\BourbakiExercisesSection

\begin{enumerate}
\def\labelenumi{\arabic{enumi}.}
\item
  设 N 为 Z-模 Z/4Z，M 为 N 的子模 2Z/4Z，且 \(j:M\to N\) 为典范单射。证明 \(\mathsf{T}(j):\mathsf{T}(\mathsf{M})\to\mathsf{T}(\mathsf{N})\) 不是单射的。
\item
  设 I 和 J 是两个有限集合， \(I_{0} = I \cup \{\alpha\}\) , \(J_{0} = J \cup \{\beta\}\) ，其中 \(\alpha\) 与 \(\beta\) 分别是不属于 I 与 J 的元素，并设 \(I_{0}\) 与 \(J_{0}\) 不相交。设 E 是交换域 K 上有限秩的向量空间，z 是 \(\mathbf{T}_{J}^{\mathrm{I}}(\mathbf{E})\) 的一个元素。对所有 \(i \in I\) ，设 \(W_{i}'\) 是 \(E^{*}\) 的由元素 \(x^{*} \in E^{*}\) （满足 \(c_{\beta}^{i}(z \otimes x^{*}) = 0\) ）组成的子空间（\(c_{\beta}^{i}\) 是指标 \(i \in I\) 与指标 \(\beta \in J_{0}\) 在 \(\mathbf{T}_{J_{0}}^{\mathrm{I}}(\mathbf{E})\) 中的缩并），并设 \(V_{i}\) 是 E 中与 \(W_{i}'\) 正交的子空间。对所有 \(j \in J\) ，设 \(W_{j}\) 是 E 的由元素 \(x \in E\) （满足 \(c_{j}^{\alpha}(z \otimes x) = 0\) ）组成的子空间（\(c_{j}^{\alpha}\) 是指标 \(\alpha \in I_{0}\) 与指标 \(j \in J\) 在 \(T_{J}^{I_{0}}(E)\) 中的缩并），并设 \(V_{j}'\) 是 \(E^{*}\) 中与 \(W_{j}\) 正交的子空间。证明 z 属于张量积 \(\left(\bigotimes_{i \in I} V_{i}\right) \otimes \left(\bigotimes_{j \in J} V_{j}'\right)\) ，它典范地等同于 \(T_{J}^{I}(E)\) 的一个子空间。此外，若 \((\mathbf{U}_{i})_{i \in I}\) 是 E 的一族子空间， \((\mathbf{U}_{j}')_{j \in J}\) 是 \(E^{*}\) 的一族子空间，使得 z 属于张量积 \(\left(\bigotimes_{i \in I} U_{i}\right) \otimes \left(\bigotimes_{j \in J} U_{j}'\right)\) ，它典范地等同于 \(T_{J}^{I}(E)\) 的一个子空间，则 \(V_{i} \subset U_{i}\) 且 \(V_{j}' \subset U_{j}'\) 对所有 \(i \in I\) 和 \(j \in J\) 成立（参见 II，§7，第 8 号，命题 17）。
\item
  设 M 是有限生成的投射 A-模。对 M 的每个自同态 u，用 \(\tilde{u}\) 表示与之对应的张量 \(\theta_{\mathbf{M}}^{-1}(u) \in \mathbf{M}^{*} \otimes_{\mathbf{A}} \mathbf{M}\) 。
\end{enumerate}

\begin{enumerate}
\def\labelenumi{(\alph{enumi})}
\item
  对每个元素 \(x \in \mathbf{M}\) ，证明 \(u(x) = c_2^1(x \otimes \tilde{u})\) ，其中 \(x \otimes \tilde{u}\) 属于 \(\mathbf{M} \otimes \mathbf{M}^* \otimes \mathbf{M} = \mathsf{T}_{\langle 2\rangle}^{(1,3)}(\mathbf{M})\) 。
\item
  设 \(u, v\) 是 \(M\) 的两个自同态且 \(w = u \circ v\) ；证明 \(\tilde{w} = c_3^2(\tilde{v} \otimes \tilde{u})\) ，其中 \(\tilde{v} \otimes \tilde{u}\) 属于 \(T_{(1,3)}^{(2,4)}(M)\) 。
\item
  对每个自同构 \(s\) （作用于 \(\mathbf{M}\)）与每个张量 \(z \in \mathbf{T}_{\mathbf{J}}^{\mathbf{I}}(\mathbf{M})\) ，张量 \(s, z \in \mathbf{T}_{\mathbf{J}}^{\mathbf{I}}(\mathbf{M})\) 由如下条件定义：若 \(z = \bigotimes_{k \in \mathbf{I} \cup \mathbf{J}} z_k\) ，其中 \(z_k \in \mathbf{M}\) （若 \(k \in \mathbf{I}\)）而 \(z_k \in \mathbf{M}^*\) （若 \(k \in \mathbf{J}\)），则 \(s, z = \bigotimes_{k \in \mathbf{I} \cup \mathbf{J}} z_k'\) ，其中 \(z_k' = s(z_k)\) （若 \(k \in \mathbf{I}\)），\(z_{k}^{\prime}=^{t}s^{-1}(z_{k})\) （若 \(k\in J\)）。证明群 \(\mathbf{GL}(\mathbf{M})\) 以法则 \((s,z)\mapsto s.z\) 作用于 A-模 \(T_{J}^{I}(M)\) 上，诸映射 \(z\mapsto s.z\) 是该 A-模的自同构。此外，对每对有序指标 \(i\in I\) , \(j\in J\) ,
\end{enumerate}

\[
c _ {j} ^ {i} (s. z) = s. c _ {j} ^ {i} (z).
\]

若 M 是投射的且有限生成的，则在习题 2 的记号下证明 \(s. \tilde{u} = (s \circ u \circ s^{-1})^{\sim}\) 。

\begin{enumerate}
\def\labelenumi{\arabic{enumi}.}
\setcounter{enumi}{3}
\tightlist
\item
  设 M 是有限生成的投射 A-模；对每个双线性映射 \(u\) （从 \(\mathbf{M} \times \mathbf{M}\) 到 \(\mathbf{M}\) 中），用 \(\tilde{u}\) 表示张量 \(\theta_{\mathbf{M}}^{-1}(u)\) （属于 \(\mathsf{T}_{(1,2)}^{(3)}(\mathbf{M})\)）。证明：\(\mathbf{M}\) 上由 \(u\) 定义的代数结构是结合的，当且仅当张量 \(c_4^3 (\tilde{u} \otimes \tilde{u})\) 与 \(c_2^6 (\tilde{u} \otimes \tilde{u})\) 在从 \(\mathsf{T}_{(1,2,3)}^{(4)}(\mathbf{M})\) 到 \(\mathsf{T}_{(1,3,4)}^{(2)}(\mathbf{M})\) 上的典范结合同构下互相对应。
\end{enumerate}

¶ 5. 设 E 是交换域 K 上的向量空间，并设 T(E) 是 E 的张量代数。

\begin{enumerate}
\def\labelenumi{(\alph{enumi})}
\item
  证明 \(\mathsf{T}(\mathbf{E})\) 不含零因子，且 \(\mathsf{T}(\mathbf{E})\) 中仅有的可逆元素是标量 \(\neq 0\) 。若 \(\dim (\mathbf{E}) > 1\) ，则 \(\mathsf{T}(\mathbf{E})\) 是非交换 K-代数。
\item
  设 u、v 是 T(E) 的两个元素。证明：若存在 T(E) 的两个元素 a、b 使得 \(ua = vb \neq 0\) ，则 u、v 中之一是另一个的右倍元。（先考虑 u 与 v 都是齐次的情形；把它化归为 E 是有限维的情形，并写出 \(u = e_{1}u_{1} + \cdots + e_{n}u_{n}\) , \(v = e_{1}v_{1} + \cdots + e_{n}v_{n}\) ，其中 \((e_{k})\) 是 E 的一个基。再对 ua 的齐次分量的最小次数作归纳得出结论。）由此推出：若 \(\dim(\mathrm{E}) > 1\) ，环 T(E) 没有左分式域（I，§ 9，习题 15）。
\item
  证明：两个非零元素 \(u, v\) （属于 \(\mathsf{T}(\mathbf{E})\)）可交换，当且仅当存在向量 \(x \neq 0\) （属于 \(\mathbf{E}\)），使得 \(u\) 与 \(v\) 属于 \(\mathsf{T}(\mathbf{E})\) 的由 \(x\) 生成的子代数（利用 (b)）。由此推出：若 \(\dim(\mathbf{E}) > 1\) ，则 \(\mathsf{T}(\mathbf{E})\) 的中心等于 \(\mathbf{K}\) 。
\end{enumerate}

¶ 6. 设 K 是交换域，E、F 是两个 K-代数，各带单位元，且该单位元与 K 的单位元 l 等同。考虑 K-代数 T(E ⊕ F)，其中 E ⊕ F 被看作 K 上的向量空间；在典范单射下，E（相应地 F）等同于 T(E ⊕ F) 的一个向量子空间。设 \(\mathfrak{J}\) 是 T(E ⊕ F) 中由元素 \(x_1 \otimes x_2 - (x_1 x_2)\) 与 \(y_1 \otimes y_2 - (y_1 y_2)\) 生成的双边理想，其中 \(x_i \in E\) , \(y_i \in F\) , \(i = 1, 2\) ；设 E * F 表示商代数 T(E ⊕ F)/\(\mathfrak{J}\)；它称为 K-代数 E 与 F 的自由积；设 \(\phi\) 是 T(E ⊕ F) 到 E * F 上的典范同态，\(\alpha\) 与 \(\beta\) 是 \(\phi\) 在 E 与 F 上的限制；\(\alpha\) 与 \(\beta\) 是 K-代数同态。

\begin{enumerate}
\def\labelenumi{(\alph{enumi})}
\item
  证明：对每一对 K-同态 \(u: \mathbf{E} \to \mathbf{G}\) , \(v: \mathbf{F} \to \mathbf{G}\) （取值在 K-代数 \(\mathbf{G}\) 中），存在唯一的 K-同态 \(w: \mathbf{E} * \mathbf{F} \to \mathbf{G}\) ，使得 \(u = w \circ \alpha\) 且 \(v = w \circ \beta\) （这就证明了该名称的合理性）。
\item
  设 \(R_{n}\) 是 \(\mathsf{T}(\mathbf{E}\oplus\mathbf{F})\) 中诸 \(\mathsf{T}^{k}(\mathbf{E}\oplus\mathbf{F})\) 对 \(0\leqslant k\leqslant n\) 之和，并记 \(J_{n}=J\cap R_{n}\) 。证明：若考虑由 l 与族 \((e_{\lambda})_{\lambda\in\mathbb{L}}\) 组成的 E 的基，以及由 l 与族 \((f_{\mu})_{\mu\in\mathbb{M}}\) 组成的 F 的基，则 \(J_{n}+R_{n-1}\) 在 \(R_{n}\) 中的一个补子空间可如下得到：在 \(R_{n}\) 中考虑以形如 \(z_{1}\otimes z_{2}\otimes\cdots\otimes z_{n}\) 的元素为基的子空间，其中或者 \(z_{2j+1}\) 等于某个 \(e_{\lambda}\) （\(2j+1\leqslant n\)）且 \(z_{2j}\) 等于某个 \(f_{\mu}\) （\(2j\leqslant n\)），或者 \(z_{2j+1}\) 等于某个 \(f_{\mu}\) （\(2j+1\leqslant n\)）且 \(z_{2j}\) 等于某个 \(e_{\lambda}\) （\(2j\leqslant n\)）（对 n 作归纳）。
\item
  设 \(\mathbf{P}_n = \phi (\mathbf{R}_n)\subset \mathbf{E}*\mathbf{F}\) ；于是 \(\mathbf{P}_n\subset \mathbf{P}_{n + 1},\mathbf{P}_0 = \mathbf{K},\mathbf{P}_h\mathbf{P}_k\subset \mathbf{P}_{h + k}\) ，且 \(\mathbf{E}*\mathbf{F}\) 是诸 \(\mathbf{P}_n\) 之并。证明存在典范的向量空间同构，对 \(n\geqslant 1\)
\end{enumerate}

\[
\mathrm{P} _ {n} / \mathrm{P} _ {n - 1} \rightarrow \left(\mathrm{G} _ {1} \otimes \mathrm{G} _ {2} \otimes \dots \otimes \mathrm{G} _ {n}\right) \oplus \left(\mathrm{G} _ {2} \otimes \mathrm{G} _ {3} \otimes \dots \otimes \mathrm{G} _ {n + 1}\right)
\]

其中记 \(G_{k} = E/K\) （k 为奇数），\(G_{k} = F/K\) （k 为偶数）。特别地，同态 \(\alpha\) 与 \(\beta\) 是单射，这使我们可以把 E 与 F 等同于 E * F 的子 K-代数。用 \(P_{n}^{\prime}\) （相应地 \(P_{n}^{\prime\prime}\) ）表示 \(G_{1} \otimes G_{2} \otimes \cdots \otimes G_{n}\) （相应地 \(G_{2} \otimes G_{3} \otimes \cdots \otimes G_{n+1}\) ）在 P 中的逆像；元素 \(z \in E * F\) 的高度（记作 \(h(z)\) ）是使 \(z \in P_{n}\) 成立的最小整数；若 \(h(z) = n\) 且 \(z \in P_{n}^{\prime}\) （相应地 \(z \in P_{n}^{\prime\prime}\) ），则称 z 为纯奇的（相应地，纯偶的）。于是 \(P_{n} = P_{n}^{\prime} + P_{n}^{\prime\prime}\) 且 \(P_{n}^{\prime} \cap P_{n}^{\prime\prime} = P_{n-1}\) 对 \(n \geqslant 1\) 成立。

\begin{enumerate}
\def\labelenumi{(\alph{enumi})}
\setcounter{enumi}{3}
\item
  证明：若 \(u, v\) 是 \(\mathbf{E} * \mathbf{F}\) 中使 \(h(u) = r\) , \(h(v) = s\) 的元素，则 \(h(uv) \leqslant h(u) + h(v)\) ，且 \(h(uv) = h(u) + h(v)\) ，除非 \(u\) 与 \(v\) 都是纯的，并且或者 \(r\) 是偶数而 \(u, v\) 奇偶性不同，或者 \(r\) 是奇数而 \(u, v\) 奇偶性相同。（例如设 \(r\) 是偶数， \(u \in \mathbf{P}_r'\) , \(v \in \mathbf{P}_s'\) ，证明 \(uv \in \mathbf{P}_{r+s-1}\) 仅当 \(u \in \mathbf{P}_{r-1}\) 或 \(v \in \mathbf{P}_{s-1}\) 时才可能。）
\item
  假设 E 没有零因子，且在 (d) 的记号下设 r 是偶数，u 是纯偶的而 v 是纯奇的。证明 \(h(uv) = r + s - 1\) ，除非存在 E 的两个可逆元素 x、y 使得 \(uy \in P_{r-1}\) 且 \(xb \in P_{s-1}\) 。（设 \((a_i)\) （相应地 \((b_j)\) ）是 (b) 中所考虑的 \(J_r + R_{r-1}\) 在 \(R_r\) 中（相应地 \(J_s + R_{s-1}\) 在 \(R_s\) 中）的补子空间的基；写出 \(u = \sum_i \phi(a_i)x_i\) 与 \(v = \sum_j y_j \phi(b_j)\) ，其中 \(x_i\) 与 \(y_j\) 都在 E 中，并证明：若 \(uv \in P_{r+s-2}\) ，则必然有 \(x_i y_j \in K\) 。）当设 E 与 F 都没有零因子时，类似地处理 \(h(uv) < h(u) + h(v)\) 的其他情形。
\item
  由(d)和(e)推出，当E和F没有零因子时，E * F也没有零因子。
\item
  将上述定义和结果推广到任意有限个含单位元的K-代数。
\end{enumerate}

\begin{enumerate}
\def\labelenumi{\arabic{enumi}.}
\setcounter{enumi}{6}
\tightlist
\item
  设 E 是域 K 上维数为 n 的向量空间。对每个整数 \(r \geqslant 1\) ，对称群 \(S_{r}\) 线性地作用于张量积 \(E^{\otimes r}\) 上，其作用为 \((\sigma, z) \mapsto \sigma \cdot z\) ，使得
\end{enumerate}

\[
\sigma . \left(x _ {1} \otimes x _ {2} \otimes \dots \otimes x _ {r}\right) = x _ {\sigma^ {- 1} (1)} \otimes x _ {\sigma^ {- 1} (2)} \otimes \dots \otimes x _ {\sigma^ {- 1} (r)}
\]

对于每个由 \(r\) 个向量组成的序列 \((x_i)_{1 \leqslant i \leqslant r}\) ，诸向量属于 \(\mathbf{E}\) 。设 \(u_{\sigma}(z) = \sigma, z\) 。证明：对每个序列 \((v_i)_{1 \leqslant i \leqslant r}\) （由 \(r\) 个 \(\mathbf{E}\) 的自同态组成），

\[
\operatorname{Tr} (u _ {\sigma} \circ (v _ {1} \otimes v _ {2} \otimes \dots \otimes v _ {r})) = \operatorname{Tr} (w _ {1}) \operatorname{Tr} (w _ {2}) \dots \operatorname{Tr} (w _ {k})\tag{*}
\]

其中 \(w_{j}\) 是 \(\mathbf{E}\) 的自同态，定义如下：把 \(\sigma^{-1}\) 分解为轮换， \(\sigma^{-1} = \lambda_1\lambda_2\ldots \lambda_k\) ，并且若 \(\lambda_{j} = (i_{1}i_{2}\dots i_{h})\) ，则记

\[
w _ {j} = v _ {i _ {1}} \circ v _ {i _ {2}} \circ \dots \circ v _ {i _ {h}}.
\]

（把它化为计算 (*) 的左端，其中每个 \(v_{i}\) 都形如 \(x \mapsto \langle x, a_{h}^{*} \rangle a_{k}\) ，这里 \((a_{j})_{1 \leqslant j \leqslant n}\) 是 E 的一个基，而 \((a_{j}^{*})_{1 \leqslant j \leqslant n}\) 是对偶基。）

¶ 8. 设 A 是一个整环，K 是其分式域；对于每个 A-模 M，令 τ(M) 表示其挠模（II，§ 7，第 10 号）。

\begin{enumerate}
\def\labelenumi{(\alph{enumi})}
\item
  证明：代数 \(\mathsf{T}(\mathbf{M})\) 没有零因子，当且仅当 A-模 \(\mathsf{T}(\mathbf{M})\) 是无挠的（注意 \(\mathsf{T}(\mathbf{M})_{(\mathbb{K})}\) 同构于 \(\mathsf{T}(\mathbf{M}_{(\mathbb{K})})\) ，并利用习题 5）。
\item
  设 \(u: \mathbf{M} \to \mathbf{M}'\) 是一个单射 A-模同态。证明 \(\operatorname{Ker}(\mathsf{T}(u)) \subset \tau(\mathsf{T}(\mathbf{M}))\) ；若 \(\mathbf{M}'\) 是无挠的，则 \(\operatorname{Ker}(\mathsf{T}(u)) = \tau(\mathsf{T}(\mathbf{M}))\) （参见 II，§ 7，第 10 号，命题 27）。
\item
  给出一个无挠 A-模 M 的例子，使得 \(\tau(\mathsf{T}(\mathbf{M}))\) 非零（II，§ 7，习题 31）。
\end{enumerate}

¶ 9. 设 A 是一个交换环，F 是整数集 I = {[}1, n{]}（其中 n ≥ 2）上的自由结合代数，而 X \(_{i}\) (1 ≤ i ≤ n) 是相应的未定元，于是 F 在 A 上的一个典范基由乘积 X \(_{i_1}\) X \(_{i_2}\) \ldots{} X \(_{i_r}\) 组成，其中 (i \(_{1}\) , \ldots， i \(_{r}\) ) 是 I 中元素的任意有限序列。

\begin{enumerate}
\def\labelenumi{(\alph{enumi})}
\item
  设 \(C_{2}\) 表示换位子 \([X_{i}, X_{j}] = X_{i}X_{j} - X_{j}X_{i}\) （\(1 \leqslant i < j \leqslant n\)）的集合。设 \(C_{m-1}\) 已有定义，\(C_{m}\) 表示元素 \([X_{i}, P] = X_{i}P - PX_{i}\) 的集合，其中 \(1 \leqslant i \leqslant n\) 而 P 取遍 \(C_{m-1}\) 。设 U 是 F 中由 l 与诸 \(C_{m}\) （\(m \geqslant 2\)）的并集生成的分次子代数。证明：对 \(1 \leqslant i \leqslant n\) ，有 \([X_{i}, P] \in U\) 对所有 \(P \in U\) （只需限于 P 是 \(\bigcup_{m} C_{m}\) 中元素之积的情形，并对因子个数作归纳论证）。
\item
  从现在起设 A 是域。对每个整数 \(m \geqslant 0\) ，设 \(\mathbf{U}_m\) 是 U 中次数为 \(m\) 的齐次元素所成的向量 A-空间（于是 \(\mathbf{U}_0 = \mathbf{A}\) , \(\mathbf{U}_1 = \{0\}\) ）；选取一组基 \(\mathbb{R}_{m,1}, \ldots, \mathbb{R}_{m,d_m}\) （在 \(\mathbf{U}_m\) 的每一个中），由 \(\bigcup_{k\leqslant m}\mathbf{C}_k\) 中元素的乘积组成。对每个有序整数对 \((m,k)\) （满足 \(m\geqslant 2,1\leqslant k\leqslant d_m\)），设 \(\mathbf{E}_{m,k}\) 是 U 中由诸 \(\mathbb{R}_{p,j}\) （满足 \(p > m\) 或 \(p = m\) 且 \(j\geqslant k\) ）生成的向量子 A-空间；证明 \(\mathbf{E}_{m,k}\) 是 U 的一个双边理想，且若 \(\mathbf{L}_{m,k}\) 是 F 中由 \(\mathbf{E}_{m,k},\mathbf{L}_{m,k}\) 生成的左理想，则它是 F 的一个双边理想。
\item
  对于每一个多重指标 \(\alpha = (\alpha_{1},\ldots ,\alpha_{n})\in \mathbf{N}^{n}\) ，我们记作
\end{enumerate}

\[
\mathrm{X} ^ {\alpha} = \mathrm{X} _ {1} ^ {\alpha_ {1}} \mathrm{X} _ {2} ^ {\alpha_ {2}} \dots \mathrm{X} _ {n} ^ {\alpha_ {n}}.
\]

证明元素 \(X^{\alpha}R_{m,k}\) ( \(m > 0, l \leqslant k \leqslant d_{m}\) ）构成 F 在 K 上的一个基。（为证明这些元素构成向量空间 F 的一个生成系，请证明每个元素 \(X_{i_{1}} \ldots X_{i_{r}}\) 都是它们的线性组合，对次数 r 作归纳。为看出 \(X^{\alpha}R_{m,k}\) 线性无关，用反证法论证：考虑其中一个 \(X_{i}\) ，并考察诸单项式 \(X^{\alpha}\) （它们出现在 \(X^{\alpha}R_{m,k}\) 之间的一个线性关系中）关于该变元的最大次数，再对此最大次数作归纳；为此，化归到 K 为无限的情形（必要时考虑 \(F \otimes_{K} K'\) ，其中 \(K'\) 是 K 的一个扩域），并注意到当 \(X_{i} + \lambda\) 代换某个 \(R_{m,k}\) 中的 \(X_{i}\) （对某个 \(\lambda \in K'\) ）时，元素 \(R_{m,k}\) 保持不变）。

\begin{enumerate}
\def\labelenumi{(\alph{enumi})}
\setcounter{enumi}{3}
\item
  证明：双边理想 \(\mathfrak{J}(\mathbf{F})\) ，即 \(\mathbf{F}\) 中由换位子 \([u,v] = uv - vu\) （u、v 为 \(\mathbf{F}\) 的元素）生成的理想，是诸 \(\mathbf{U}_m\) （\(m\geqslant 2\)）之和；\(\mathbf{F} / \mathfrak{J}(\mathbf{F})\) 同构于 \(\mathrm{A}[\mathrm{X}_1,\dots ,\mathrm{X}_n]\) ，从而是一个整环。
\item
  证明代数 U 不是自由结合代数，方法是证明商 \(\mathrm{U}/\Im(\mathrm{U})\) （其中 \(\Im(\mathrm{U})\) 由 U 中元素的换位子生成）含有幂零元素 \(\neq0\) 。（考虑 \(z'\) ，即 \(\mathrm{U}/\Im(\mathrm{U})\) 中元素 \(z = [\mathbf{X}_{i}, \mathbf{X}_{j}]\) 的像；\(z'^{2} \neq 0\) 但 \(z'^{4} = 0\) 。）
\end{enumerate}

\begin{enumerate}
\def\labelenumi{\arabic{enumi}.}
\setcounter{enumi}{9}
\tightlist
\item
  设 A 为一个交换环，
\end{enumerate}

\[
\mathrm{M} ^ {\prime} \xrightarrow {u} \mathrm{M} \xrightarrow {v} \mathrm{M} ^ {\prime \prime} \longrightarrow 0
\]

一个 A-模的正合序列，并设 \(n\) 为一个整数 \(>0\) 。设 \(L = \operatorname{Coker}(T^n(u))\) ，于是有正合序列

\[
\mathsf {T} ^ {n} (\mathbf {M} ^ {\prime}) \xrightarrow {\mathsf {T} ^ {n} (u)} \mathsf {T} ^ {n} (\mathbf {M}) \to \mathbf {L} \to 0.
\]

对集合 \(\{1, 2, \ldots, n\}\) 的每个子集 J ，我们记

\[
\mathrm{P} _ {J} = \mathrm{N} _ {1} \otimes \mathrm{N} _ {2} \otimes \dots \otimes \mathrm{N} _ {n}, \quad \text { where } \mathrm{N} _ {i} = \mathrm{M} ^ {\prime} \text { if } i \in \mathrm{J}, \mathrm{N} _ {i} = \mathrm{M} \text { if } i \notin \mathrm{J}
\]

\((\mathbf{P}_{\varnothing}\) 因此等于 \(\mathsf{T}^n (\mathbf{M}))\) 和

\[
\mathrm{Q} _ {\mathrm{J}} = \mathrm{N} _ {1} \otimes \mathrm{N} _ {2} \otimes \dots \otimes \mathrm{N} _ {n}, \quad \text { where } \mathrm{N} _ {i} = \mathrm{M} ^ {\prime} \text { if } i \in \mathrm{J}, \mathrm{N} _ {i} = \mathrm{M} ^ {\prime \prime} \text { if } i \notin \mathrm{J}
\]

\((Q_{\varnothing}\) 因此等于 \(T^{n}(M^{\prime \prime})\) 。）对每个整数 \(i\) （满足 \(0 \leqslant i \leqslant n\)），设 \(P^{(i)}\) 表示 \(T^{n}(M)\) 的子模，它由诸 \(\mathbf{P}_{\mathbf{J}}\) （\(\mathbf{J}\) 取遍满足 \(\operatorname{Card}(\mathbf{J}) = i\) 的一切子集）的像的并集生成；并设 \(\mathbf{L}^{(i)}\) 为 \(\mathbf{P}^{(i)}\) 在 \(\mathbf{L}\) 中的典范像。

\begin{enumerate}
\def\labelenumi{(\alph{enumi})}
\tightlist
\item
  证明存在一个典范的满同态
\end{enumerate}

\[
\bigoplus_ {\operatorname{Card} (J) = i} Q _ {J} \rightarrow L ^ {(i)} / L ^ {(i + 1)}
\]

(参见 II, § 5, 第 6 号, 命题 6)。

\begin{enumerate}
\def\labelenumi{(\alph{enumi})}
\setcounter{enumi}{1}
\tightlist
\item
  若序列 \(0 \to \mathbf{M}' \xrightarrow{u} \mathbf{M} \xrightarrow{v} \mathbf{M}'' \to 0\) 是正合的且分裂，证明在 (a) 中定义的同态是双射。
\end{enumerate}

\BourbakiExercisesSection

\begin{enumerate}
\def\labelenumi{\arabic{enumi}.}
\item
  若 M 是单生成 A-模，则 \(\mathsf{S}(\mathsf{M}) = \mathsf{T}(\mathsf{M})\) 。由此推出一个单同态 \(j: N \to M\) 的例子，使得 \(\mathsf{S}(j): \mathsf{S}(\mathsf{N}) \to \mathsf{S}(\mathsf{M})\) 不是单射（ \(\S 5\) ，习题 1）。
\item
  为对称代数建立与§5习题8(a)和(b)中性质相类似的性质。
\item
  设 \(\mathbf{K}\) 是一个域， \(\mathbf{B}\) 为多项式环 \(\mathbf{K}[\mathbf{X}_1, \mathbf{X}_2]\) ，而 \(\mathfrak{p}\) 为 \(\mathbf{B}\) 中由多项式 \(\mathbf{X}_1^3 - \mathbf{X}_2^2\) 生成的理想。
\end{enumerate}

\begin{enumerate}
\def\labelenumi{(\alph{enumi})}
\item
  证明 p 是素理想，因此 A = B/p 是整环。设 a 和 b 分别表示 \(X_{1}\) 和 \(X_{2}\) 在 A 中的典范像；a 和 b 在 A 中不可逆，且 \(a^{3} = b^{2}\) 。
\item
  设 m 是 A 的由 a 和 b 生成的（极大）理想。证明对称代数 \(\mathsf{S}(\mathsf{m})\) 不是整环。（注意，m 与商模 \(A^{2}/N\) 等同，其中 N 是 \(A^{2}\) 的由 \((b, -a)\) 生成的子模；得出结论：\(\mathsf{S}(\mathsf{m})\) 同构于商环 \(\mathbf{A}[\mathbf{U}, \mathbf{V}]/(b\mathbf{U} - a\mathbf{V})\) ，并证明在多项式环 \(\mathbf{A}[\mathbf{U}, \mathbf{V}]\) 中主理想 \((b\mathbf{V} - a\mathbf{U})\) 不是素理想，这可通过考虑乘积 \(b^{2}(b\mathbf{V}^{3} - \mathbf{U}^{3})\) 来实现。）
\end{enumerate}

¶ 4. 设 A 是一个交换环，a 是 A 的一个理想。理想 a 的里斯代数是单未定元多项式代数 A{[}T{]} 的由形如 \(c_{0} + c_{1}\mathrm{T} + \cdots + c_{n}\mathrm{T}^{n}\) 的多项式组成的子代数 R(a)，其中 \(c_{0} \in \mathbf{A}\) ，且 \(c_{k} \in \mathfrak{a}^{k}\) （对每个整数 \(k \geqslant 1\)）。

\begin{enumerate}
\def\labelenumi{(\alph{enumi})}
\item
  证明存在且仅存在一个满射 A-代数同态 \(r: \mathsf{S}(\mathfrak{a}) \to \mathsf{R}(\mathfrak{a})\) ，使得 \(r\) 在 \(\mathfrak{a}\) 上的限制是 A-线性映射 \(x \to x\mathbb{T}\) 。
\item
  设环 A 是整环。令 \(a_1, \ldots, a_n\) 是 A 的元素，且 \(a_n \neq 0\) ；考虑 A-代数同态
\end{enumerate}

\[
u: \mathrm{A} [ \mathrm{X} _ {1}, \dots , \mathrm{X} _ {n} ] \rightarrow \mathrm{A} [ \mathrm{T} ]
\]

，使得 \(u(\mathbf{X}_i) = a_i\mathrm{T}\) 。核 \(\mathfrak{r} = \operatorname{Ker}(u)\) 是一个分次理想，其齐次分量 \(\mathfrak{r}_m\) （次数为 \(m\) ）由齐次多项式 \(f\) （次数为 \(m\) 且满足 \(f(a_1, \ldots, a_n) = 0\) ）组成。若 \(\mathfrak{r}' \subset \mathfrak{r}\) 是 \(\mathrm{A}[\mathrm{X}_1, \ldots, \mathrm{X}_n]\) 中由齐次分量 \(\mathfrak{r}_1\) （属于 \(\mathfrak{r}\)）生成的理想，证明 A-模 \(\mathfrak{r}/\mathfrak{r}'\) 是一个挠 A-模。（为验证：对所有 \(f \in \mathfrak{r}_m\) ，存在 A 中的 \(c \neq 0\) 使得 \(cf \in \mathfrak{r}'\) ，对 \(m\) 作归纳论证，写出

\[
f \left(\mathrm{X} _ {1}, \dots , \mathrm{X} _ {n}\right) = \mathrm{X} _ {1} f _ {1} \left(\mathrm{X} _ {1}, \dots , \mathrm{X} _ {n}\right) + \mathrm{X} _ {2} f _ {2} \left(\mathrm{X} _ {2}, \dots , \mathrm{X} _ {n}\right) + \dots + \mathrm{X} _ {n} f _ {n} \left(\mathrm{X} _ {n}\right)
\]

其中 \(f_{j}\) 是次数 \(m - 1\) 的齐次多项式；考虑多项式 \(g(\mathbf{X}_1, \ldots, \mathbf{X}_n) = \mathbf{X}_1 f_1(a_1, \ldots, a_n) + \mathbf{X}_2 f_2(a_2, \ldots, a_n) + \cdots + \mathbf{X}_n f_n(a_n)\) ，观察到 \(g \in \mathfrak{r}_1 \subset \mathfrak{r}'\) ，并作差 \(a_n^{m-1} f - \mathbf{X}_n^{m-1} g\) 。

\begin{enumerate}
\def\labelenumi{(\alph{enumi})}
\setcounter{enumi}{2}
\tightlist
\item
  从 (b) 推断，当 A 是整环时，(a) 中同态 r 的核是 S(a) 的挠模。由此推断，为使 S(a) 是整环，必须且只需 S(a) 是无挠 A-模。特别地，若理想 a 是投射 A-模，则 S(a) 是整环且同构于 R(a)。
\end{enumerate}

\begin{enumerate}
\def\labelenumi{\arabic{enumi}.}
\setcounter{enumi}{4}
\tightlist
\item
  设 E 是域 K 上维数有限为 n 的向量空间。
\end{enumerate}

\begin{enumerate}
\def\labelenumi{(\alph{enumi})}
\tightlist
\item
  证明对每个整数 \(m\) ，对称幂 \(S^m(E)\) 与向量空间 \(S_m'(E)\) （即 \(m\) 阶对称逆变张量的空间）具有相同的维数 \(\binom{m + n - 1}{n - 1} = \binom{n + m - 1}{m}\) 。
\end{enumerate}

*(b) 若 K 的特征为 p \textgreater{} 0，证明 p 阶对称化张量所成的向量空间 \(\mathbb{S}_{p}^{\prime\prime}(\mathbf{E})\) 的维数为 \(\binom{n}{p}\) （注意，若在由 E 的 p 个向量组成的张量积 \(x_{1} \otimes x_{2} \otimes \cdots \otimes x_{p}\) 中，对一对不同的指标 i, j 有 \(x_{i} = x_{j}\) ，则该乘积的对称化为零；利用如下事实：若 \((\mathrm{H}_{i})_{1 \leqslant i \leqslant r}\) 是 \(\{1, 2, \ldots, p\}\) 的一个分成非空集合（至少含两个集合）的划分，则 \(S_{p}\) 中保持每个 \(H_{i}\) 稳定的子群的阶不被 p 整除）。另一方面，在 \(\mathbb{S}^{p}(\mathbf{E})\) 中，p 阶对称张量空间 \(\mathbb{S}_{p}^{\prime}(\mathbf{E})\) 的典范像的维数为 n，且由张量 \(x \otimes x \otimes \cdots \otimes x\) （p 次）的像生成，其中 \(x \in E\) （利用同一注记）；因此从 \(\mathbb{S}_{p}^{\prime}(\mathbf{E})\) 到 \(\mathbb{S}^{p}(\mathbf{E})\) 的典范映射不是双射，尽管这两个向量空间具有相同的维数。*

\BourbakiExercisesSection

\begin{enumerate}
\def\labelenumi{\arabic{enumi}.}
\tightlist
\item
  设 A 为整环，K 为其分式域。
\end{enumerate}

\begin{enumerate}
\def\labelenumi{(\alph{enumi})}
\item
  证明外幂 \(\bigwedge^{2}\mathbf{K}\) （即 A-模 \(\mathbf{K}\) 的外幂）缩减为 0。
\item
  对每个 A-模 \(\mathbf{E} \subset \mathbf{K}\) ，证明每个外幂 \(\bigwedge \mathbf{E}\) （\(p \geqslant 2\)）都是挠 A-模。
\item
  若 E 为 II，§7，习题 31 中定义的 A-模，证明 \(\bigwedge\) E 是非零的。给出一个整环 A 和包含在 A 的分式域 K 中的 A-模 E 的例子，使得任何外幂 \(\bigwedge^{p}\) E 都不化为 0。
\end{enumerate}

\begin{enumerate}
\def\labelenumi{\arabic{enumi}.}
\setcounter{enumi}{1}
\item
  对于环 A 的每个理想 a 和每个 A-模 E，外代数 \(\bigwedge(\mathrm{E}/\mathfrak{a}\mathrm{E})\) 同构于 \((\bigwedge\mathrm{E})/\mathfrak{a}(\bigwedge\mathrm{E})\) 。
\item
  给出一个具有如下性质的环 A 的例子：在 A-模 \(E = A^{2}\) 中，存在一个子模 F，使得，若 \(j: F \to E\) 是典范单射，则 \(\bigwedge^{2} j\) 恒为零，即使 \(\bigwedge^{2} E\) 与 \(\bigwedge^{2} F\) 都不化为 0（II，§4，习题 2）。
\item
  设 E 是维数为 \(n \geqslant 3\) 的自由 A-模。证明：若 \(l \leqslant p \leqslant n - 1\) ，则 A-模 \(\bigwedge^{p}E\) 与 \(\bigwedge^{n-p}E\) 同构；但若 \(2p \neq n\) ，则不存在同构 \(\phi\) （从 \(\bigwedge^{p}E\) 到 \(\bigwedge^{n-p}E\) 上），使得对 E 的每个自同构 u 都有 \(\phi \circ \left(\bigwedge^{p}u\right) = \left(\bigwedge^{n-p}u\right) \circ \phi\) （若 \((e_{i})\) 是 E 的一组基，考虑满足 \(u(e_{i}) = e_{i} + e_{j}, y(e_{k}) = e_{k}\) （对 \(k \neq i\) ）的自同构 u，并让 i 和 j 取遍所有可能的值）。
\item
  设 K 是一个交换域，E、F 是 K 上的两个向量空间，u 是从 E 到 F 的秩为 r 的线性映射。证明：若 \(p \leqslant r\) ，则 \(\bigwedge^{p}u\) 的秩等于 \(\binom{r}{p}\) ；且若 p \textgreater{} r，则 \(\bigwedge^{p}u\) 恒为零（在 E 中取一组基，使其中 r 个向量构成与 \(u^{-1}(0)\) 互补的一个子空间的基，而其余向量构成 \(u^{-1}(0)\) 的基）。
\end{enumerate}

¶ 6. 设 K 为一个交换域，E 为 K 上的向量空间。

\begin{enumerate}
\def\labelenumi{(\alph{enumi})}
\item
  为使元素 \(z = \sum_{p=0}^{\infty} z_p \left( \text{其中 } z_p \in \bigwedge^p \mathbf{E} \right)\) （属于外代数 \(\bigwedge \mathbf{E}\) ，即 \(\mathbf{E}\) 的外代数）可逆，必须且只需 \(z_0 \neq 0\) （先在 \(\mathbf{E}\) 为有限维时证明，然后过渡到一般情形，并注意到对所有 \(z \in \bigwedge \mathbf{E}\) ，存在一个有限维子空间 \(\mathbf{F}\) （包含于 \(\mathbf{E}\)），使得 \(z \in \bigwedge \mathbf{F}\) ）。
\item
  设 K 满足在 K 中 \(2 \neq 0\) 。证明：若 E 是无限维的或有限维且维数为偶数，则代数 \(\Lambda E\) 的中心由这样的元素组成：对所有奇数指标 p，\(z_{p} = 0\) 。若 E 的维数 n 为奇数，则 \(\Lambda E\) 的中心是上述子空间与 \(\Lambda^{n}E\) 之和。在所有情形下， \(\Lambda E\) 的中心都与 \(\Lambda E\) 中与 \(\Lambda E\)
\end{enumerate}

\begin{enumerate}
\def\labelenumi{\arabic{enumi}.}
\setcounter{enumi}{6}
\item
  的所有可逆元素可交换的元素所成的集合相同。对每个 A-模 E，证明：若我们记 \([u, v] = u \wedge v - v \wedge u\) （对 \(\bigwedge E\) 的任意两个元素），则对任意三个元素 \(u, v, w\) （属于 \(\bigwedge E\)），有 \([(u, v], w] = 0\) 。
\item
  设 E 是一个自由 A-模。
\end{enumerate}

\begin{enumerate}
\def\labelenumi{(\alph{enumi})}
\tightlist
\item
  证明在 \(\mathbf{T}^n (\mathbf{E})\) 中， \(\mathfrak{J}_n''\) （第 1 号）是一个自由子 A-模，等于自同态 \(z\mapsto a.z\) 的核，并且有一个补，该补是自由 A-
\end{enumerate}

模；因此 \(\bigwedge E\) 典范同构于 \(\mathsf{A}_n^{\prime \prime}(\mathbf{E})\) 。若此外在 A 中方程 \(2\xi = 0\) 蕴含 \(\xi = 0\) ，则 \(\mathsf{A}_n^{\prime \prime}(\mathbf{E}) = \mathsf{A}_n^{\prime}(\mathbf{E})\) 。

*(b) 设 A 是特征为 p \textgreater{} 0 的域，且 \(p \neq 2\) 。则 \(\mathbf{A}_{p}^{\prime}(\mathbf{E}) = \mathbf{A}_{p}^{\prime\prime}(\mathbf{E})\) 在 \(\bigwedge^{p} E\) 中的典范像为零。

\begin{enumerate}
\def\labelenumi{(\alph{enumi})}
\setcounter{enumi}{2}
\tightlist
\item
  设 \(\mathbf{A}\) 是特征为 2 的域。则 \(\mathbf{A}_n^\prime (\mathbf{E}) = \mathbf{S}_n^\prime (\mathbf{E}),\) \(\mathbf{A}_n''(\mathbf{E}) = \mathbf{S}_n''(\mathbf{E})\) ，但一般地 \(\mathbf{A}_n^\prime (\mathbf{E})\neq \mathbf{A}_n''(\mathbf{E})\cdot *\)
\end{enumerate}

\begin{enumerate}
\def\labelenumi{\arabic{enumi}.}
\setcounter{enumi}{8}
\tightlist
\item
  设 E 是一个 A-模。从 \(\mathsf{T}^{n}(\mathbf{E})\) 到一个 A-模 F 中的线性映射 g 的斜对称化，按定义是从 \(\mathsf{T}^{n}(\mathbf{E})\) 到 F 中的线性映射 ag，它使得 \(ag(z) = g(\boldsymbol{a}z)\) 对所有 \(z \in \mathsf{T}^{n}(\mathbf{E})\) 成立；若 \(\phi: E^{n} \to T^{n}(\mathbf{E})\) 是典范映射，则 n-线性映射 \(ag \circ \phi\) 也称为 n-线性映射 \(g \circ \phi\) 的斜对称化。由习题 8(a) 推出，若 E 是自由的，则每个从 \(E^{n}\) 到 F 中的交错 n-线性映射都是某个从 \(E^{n}\) 到 F 中的 n-线性映射的斜对称化。
\end{enumerate}

¶ 10. 设 E 是一个具有 n 个元素的生成系的 A-模。证明每个从 \(E^{n}\) 到一个 A-模 F 中的交错 n-线性映射都是某个从 \(E^{n}\) 到 F 中的 n-线性映射的斜对称化。（设 \(a_{j}\) ( \(1 \leqslant j \leqslant n\) ) 为 E 的生成元，且 \(\phi: A^{n} \rightarrow E\) 为满足 \(\phi(e_{j}) = a_{j}\) （对所有 j）的同态，其中 \((e_{j})\) 是 \(A^{n}\) 的典范基。若 \(f: E^{n} \rightarrow F\) 是交错 n-线性映射，则 \(f_{0}: (x_{1}, \ldots, x_{n}) \mapsto f(\phi(x_{1}), \ldots, \phi(x_{n}))\) 也是；证明：若 \(g_{0}\) 是从 \((A^{n})^{n}\) 到 F 中的 n-线性映射，使得 \(f_{0} = a g_{0}\) ，则 \(g_{0}\) 也可以写成 \((x_{1}, \ldots, x_{n}) \mapsto g(\phi(x_{1}), \ldots, \phi(x_{n}))\) ，其中 g 是从 \(E^{n}\) 到 F 中的 n-线性映射。）

¶ 11. 设 K 是一个交换域，其中 \(2 = 0\) ，A 为多项式环 \(\mathrm{K}[\mathrm{X}, \mathrm{Y}, \mathrm{Z}]\) ，E 为 A 模 \(\mathrm{A}^3/\mathrm{M}\) ，其中 M 是 \(\mathrm{A}^3\) 的由元素 \((\mathrm{X}, \mathrm{Y}, \mathrm{Z})\) （属于 \(\mathrm{A}^3\)）生成的子模；最后，设 F 为 A 关于理想 \(\mathrm{AX}^2 + \mathrm{AY}^2 + \mathrm{AZ}^2\) 的商 A-模。证明存在一个从 \(\mathrm{E}^2\) 到 F 中的交错双线性映射，它不是任何从 \(\mathrm{E}^2\) 到 F 中的双线性映射的斜对称化，但在 \(\mathsf{T}^2(\mathbf{E})\) 中模 \(\mathfrak{S}_2''\) 等于自同态 \(z \mapsto az\) 的核（把 \(\mathsf{T}^2(\mathbf{E})\) 看作 \(\mathsf{T}^2 (\mathbf{A}^3)\) 的商模，并证明 \(z\mapsto az\) 的核是对称张量模 \(\mathsf{S}_2^{\prime}(\mathbf{A}^3)\) 的典范像）。

¶ 12. 在习题11的记号下，设B为商环

\[
\mathrm{A} / (\mathrm{AX} ^ {2} + \mathrm{AY} ^ {2} + \mathrm{AZ} ^ {2})
\]

并设 G 为 B-模 E ⊗A B。证明在 T²(G) 中，模 \(J_{2}^{\prime\prime}\) 不同于自同态 \(z \mapsto az\) 的核（方法类似于习题 11 的方法）。

\begin{enumerate}
\def\labelenumi{\arabic{enumi}.}
\setcounter{enumi}{12}
\tightlist
\item
  \begin{enumerate}
  \def\labelenumii{(\alph{enumii})}
  \tightlist
  \item
    设 M 是 N 型分次 A-模。证明在代数 \(S(M)\) （相应地 \(\bigwedge(M)\)）上存在且仅存在一种 N 型分次，使得 \(S(M)\) （相应地 \(\bigwedge(M)\)）成为分次 A-代数，并在 M 上诱导给定的分次。证明：若 M 中偶数次齐次元素全为零（此时 M 上的分次称为奇分次），则如此得到的分次代数 \(\bigwedge(M)\) 是交错的（ \(\S 4\) ，第 9 号，定义 7）。
  \end{enumerate}
\end{enumerate}

\begin{enumerate}
\def\labelenumi{(\alph{enumi})}
\setcounter{enumi}{1}
\tightlist
\item
  给定一个 \(\mathbf{N}\) 型分次 A-模，考虑以下泛映射问题： \(\Sigma\) 是反交换分次代数结构（§ 4，第 9 号，定义 7）的种，态射是分次代数的 A-同态（§ 3，第 1 号），而 \(\alpha\) -映射是 M 到反交换分次 A-代数中的次数为 0 的分次同态。设 \(\mathbf{M}^{-}\) （相应地 \(\mathbf{M}^{+}\)）表示 M 的由奇（相应地，偶）次数齐次元素生成的分次子模。证明分次代数 \(\mathbf{G}(\mathbf{M}) = \bigwedge (\mathbf{M}^{-}) \otimes_{\mathbf{A}} \mathbf{S}(\mathbf{M}^{+})\) 是反交换的；M 到这个代数中的一个典范单射 \(j\) 由写出 \(j(x) = x \otimes 1\) 若 \(x \in \mathbf{M}^{-}, j(x) = 1 \otimes x\) 若 \(x \in \mathbf{M}^{+}\) 来定义。证明分次代数 \(\mathbf{G}(\mathbf{M})\) 与单射 \(j\) 构成上述泛映射问题的一个解。
\end{enumerate}

\[
\mathbf {G} (\mathbf {M}) = \bigwedge (\mathbf {M} ^ {-}) \otimes_ {\mathbf {A}} \mathbf {S} (\mathbf {M} ^ {+})
\]

称为分次A-模M上的泛反交换代数。

\begin{enumerate}
\def\labelenumi{(\alph{enumi})}
\setcounter{enumi}{2}
\tightlist
\item
  证明关于泛反交换代数的命题 2、3 和 4 的类似结论。考虑 M 允许有由齐次元素组成的一组基的情形：在此情形下，明确给出泛反交换代数的一组基。
\end{enumerate}

\begin{enumerate}
\def\labelenumi{\arabic{enumi}.}
\setcounter{enumi}{13}
\tightlist
\item
  设 M 是一个投射 A-模，且令 \(x_{1}, \ldots, x_{n}\) 是 M 的线性无关元素；对每个子集 \(H = \{i_{1}, \ldots, i_{m}\}\) ——由 \(m \leqslant n\) 个区间 \([1, n]\) 的元素组成——其中 \(i_{1} < i_{2} < \cdots < i_{m}\) ，我们记
\end{enumerate}

\[
x _ {\mathrm{H}} = x _ {i _ {1}} \wedge x _ {i _ {2}} \wedge \dots \wedge x _ {i _ {m}}.
\]

证明当 H 遍历 \([1, n]\) 的具有 m 个元素的子集时，诸 \(x_{H}\) 在 \(\bigwedge(M)\) 中线性无关。

\begin{enumerate}
\def\labelenumi{\arabic{enumi}.}
\setcounter{enumi}{14}
\tightlist
\item
  设 M 是一个 n 维自由 A-模。证明 M 的每个具有 n 个元素的生成系都是 M 的一组基。
\end{enumerate}

\BourbakiExercisesSection

\begin{enumerate}
\def\labelenumi{\arabic{enumi}.}
\item
  在一个 n 阶矩阵 U 中，如果对于每个指标 i，将指标 i 的列替换为指标 \(\neq i\) 的列之和，则所得矩阵的行列式等于 \((-1)^{n-1}(n-1)\det U\) 。如果在 U 中，对于每个 i，从指标 i 的列中减去指标 \(\neq i\) 的列之和，则所得的行列式等于 \(-(n-2)2^{n-1}\det U\) 。
\item
  设 \(\Delta = \det (\alpha_{ij})\) 是一个 \(n\) 阶矩阵的行列式；对于
\end{enumerate}

\[
1 \leqslant i \leqslant n - 1
\]

和 \(1 \leqslant j \leqslant n - 1\) ，我们记作

\[
\beta_ {i j} = \left| \begin{array}{c c} \alpha_ {\mathbf {1} j} & \alpha_ {\mathbf {1}, j + \mathbf {1}} \\ \alpha_ {i + \mathbf {1}, j} & \alpha_ {i + \mathbf {1}, j + \mathbf {1}} \end{array} \right|.
\]

证明 n - 1 阶行列式 \(\det(\beta_{ij})\) 等于

\[
\alpha_ {1 2} \alpha_ {1 3} \dots \alpha_ {1, n - 1} \Delta .
\]

\begin{enumerate}
\def\labelenumi{\arabic{enumi}.}
\setcounter{enumi}{2}
\tightlist
\item
  证明恒等式
\end{enumerate}

\[
\left| \begin{array}{cccccc} x _ {1} & y _ {1} & \alpha_ {1 3} & \alpha_ {1 4} & \ldots & \alpha_ {1 n} \\ \lambda_ {1} x _ {2} & x _ {2} & y _ {2} & \alpha_ {2 4} & \ldots & \alpha_ {2 n} \\ \lambda_ {1} \lambda_ {2} x _ {3} & \lambda_ {2} x _ {3} & x _ {3} & y _ {3} & \ldots & \alpha_ {3 n} \\ \cdot & \cdot & \cdot & \cdot & \cdot & \cdot \\ \lambda_ {1} \lambda_ {2} \ldots \lambda_ {n - 1} x _ {n} & \lambda_ {2} \ldots \lambda_ {n - 1} x _ {n} & \lambda_ {3} \ldots \lambda_ {n - 1} x _ {n} & \lambda_ {4} \ldots \lambda_ {n - 1} x _ {n} & \ldots & x _ {n} \\ \end{array} \right| = (x _ {1} - \lambda_ {1} y _ {1}) (x _ {2} - \lambda_ {2} y _ {2}) \ldots (x _ {n - 1} - \lambda_ {n - 1} y _ {n - 1}) x _ {n}.
\]

推导出以下恒等式：

\[
\left| \begin{array}{c c c c c} 1 & 1 & 1 & \dots & 1 \\ b _ {1} & a _ {1} & a _ {1} & \dots & a _ {1} \\ b _ {1} & b _ {2} & a _ {2} & \dots & a _ {2} \\ \cdot & \cdot & \cdot & \cdot & \cdot \\ b _ {1} & b _ {2} & b _ {3} & \dots & a _ {n} \end{array} \right| = (a _ {1} - b _ {1}) (a _ {2} - b _ {2}) \dots (a _ {n} - b _ {n})
\]

\[
\left| \begin{array}{c c c c c} a _ {1} b _ {1} & a _ {1} b _ {2} & a _ {1} b _ {3} & \ldots & a _ {1} b _ {n} \\ a _ {1} b _ {2} & a _ {2} b _ {2} & a _ {2} b _ {3} & \ldots & a _ {2} b _ {n} \\ a _ {1} b _ {3} & a _ {2} b _ {3} & a _ {3} b _ {3} & \ldots & a _ {3} b _ {n} \\ \cdot \cdot & \cdot \cdot & \cdot \cdot & \cdot \cdot & \cdot \cdot \\ a _ {1} b _ {n} & a _ {2} b _ {n} & a _ {3} b _ {n} & \ldots & a _ {n} b _ {n} \end{array} \right| = a _ {1} b _ {n} (a _ {2} b _ {1} - a _ {1} b _ {2}) \ldots (a _ {n} b _ {n - 1} - a _ {n - 1} b _ {n})
\]

\[
\left| \begin{array}{cccccc} a _ {2} a _ {3} \ldots a _ {n} & a _ {3} a _ {4} \ldots a _ {n} b _ {1} & a _ {4} \ldots a _ {n} b _ {1} b _ {2} & \ldots & b _ {1} b _ {2} b _ {3} \ldots b _ {n - 1} \\ b _ {2} b _ {3} \ldots b _ {n} & a _ {3} a _ {4} \ldots a _ {n} a _ {1} & a _ {4} \ldots a _ {n} a _ {1} b _ {2} & \ldots & a _ {1} b _ {2} b _ {3} \ldots b _ {n - 1} \\ a _ {2} b _ {3} \ldots b _ {n} & b _ {3} b _ {4} \ldots b _ {n} b _ {1} & a _ {4} \ldots a _ {n} a _ {1} a _ {2} & \ldots & a _ {1} a _ {2} b _ {3} \ldots b _ {n - 1} \\ \cdot & \cdot & \cdot & \cdot & \cdot \\ a _ {2} a _ {3} \ldots b _ {n} & a _ {3} a _ {4} \ldots b _ {n} b _ {1} & a _ {4} \ldots b _ {n} b _ {1} b _ {2} & \ldots & a _ {1} a _ {2} a _ {3} \ldots a _ {n - 1} \\ \end{array} \right| = (a _ {1} a _ {2} \ldots a _ {n} - b _ {1} b _ {2} \ldots b _ {n}) ^ {n - 1}.
\]

\[
\left| \begin{array}{c c c c c c} x & a _ {1} & a _ {2} & \ldots & a _ {n - 1} & 1 \\ a _ {1} & x & a _ {2} & \ldots & a _ {n - 1} & 1 \\ a _ {1} & a _ {2} & x & \ldots & a _ {n - 1} & 1 \\ . & . & . & . & . & . \\ a _ {1} & a _ {2} & a _ {3} & \ldots & x & 1 \\ a _ {1} & a _ {2} & a _ {3} & \ldots & a _ {n} & 1 \end{array} \right| = (x - a _ {1}) (x - a _ {2}) \ldots (x - a _ {n})
\]

\[
\begin{array}{r l} & {\left| \begin{array}{l l l l l} x & a _ {1} & a _ {2} & \ldots & a _ {n} \\ a _ {1} & x & a _ {2} & \ldots & a _ {n} \\ a _ {1} & a _ {2} & x & \ldots & a _ {n} \\ \cdot & \cdot & \cdot & \cdot & \cdot \\ a _ {1} & a _ {2} & a _ {3} & \ldots & x \end{array} \right|} \\ & {\qquad = (x + a _ {1} + a _ {2} + \dots + a _ {n}) (x - a _ {1}) (x - a _ {2}) \ldots (x - a _ {n})} \end{array}
\]

（将最后一个行列式化简为前一个行列式）。

\begin{enumerate}
\def\labelenumi{\arabic{enumi}.}
\setcounter{enumi}{3}
\tightlist
\item
  计算行列式
\end{enumerate}

\[
\Delta_ {n} = \left| \begin{array}{c c c c c} a _ {1} + b _ {1} & b _ {1} & b _ {1} & \ldots & b _ {1} \\ b _ {2} & a _ {2} + b _ {2} & b _ {2} & \ldots & b _ {2} \\ \cdot & \cdot & \cdot & \cdot & \cdot \\ b _ {n} & b _ {n} & b _ {n} & \ldots & a _ {n} + b _ {n} \end{array} \right|
\]

（将 \(\Delta_{n}\) 用 \(\Delta_{n-1}\) 表出）。

\begin{enumerate}
\def\labelenumi{\arabic{enumi}.}
\setcounter{enumi}{4}
\tightlist
\item
  若 \(a_i\) 与 \(b_j\) 是交换域的元素，使得 \(a_i + b_j \neq 0\) （对每一对有序指标 \((i,j)\)），证明
\end{enumerate}

\[
\det \left(\frac {1}{a _ {i} + b _ {j}}\right) = \frac {\sum_ {i <   j} (a _ {j} - a _ {i}) (b _ {j} - b _ {i})}{\sum_ {i , j} (a _ {i} + b _ {j})}
\]

（“柯西行列式”）。

\begin{enumerate}
\def\labelenumi{\arabic{enumi}.}
\setcounter{enumi}{5}
\tightlist
\item
  证明：若 X 是 n 行 m 列的矩阵，Y 是 p 行 n 列的矩阵，且 Z = YX 是 p 行 m 列的乘积矩阵，则 Z 的 q 阶子式在 n \textless{} q 时为零；若 \(q \leqslant n\) ，它们由
\end{enumerate}

\[
\det (Z _ {L, H}) = \sum_ {K} \det (Y _ {L, K}) \det (X _ {K, H})
\]

其中 K 遍历 \([1, n]\) 的具有 q 个元素的子集所成的集合（利用 §7 第 2 号中的公式 (3)）。

\begin{enumerate}
\def\labelenumi{\arabic{enumi}.}
\setcounter{enumi}{6}
\item
  设 \(\Delta = \det(\alpha_{ij})\) 是 n 阶矩阵 U 的行列式；对每个指标 i（ \(1 \leqslant i \leqslant n\) ），令 \(\Delta_i\) 表示将 U 中的元素 \(\alpha_{ij}\) 乘以 \(\beta_j\) （对 \(1 \leqslant j \leqslant n\)）后所得的行列式。证明这 n 个行列式 \(\Delta_i\) ( \(1 \leqslant i \leqslant n\) ）之和等于 \((\beta_1 + \beta_2 + \cdots + \beta_n)\Delta\) （按指标 i 的行展开 \(\Delta_i\) ）。
\item
  设 \(\Delta = \det(\alpha_{ij})\) 是矩阵 \(U\) （\(n\) 阶）的行列式， \(\sigma\) 是属于 \(\mathfrak{S}_n\) 的一个置换；对每个指标 \(i\) ( \(1 \leqslant i \leqslant n\) )，令 \(\Delta_i\) 为把元素 \(\alpha_{ij}\) 在 \(U\) 中换成 \(\alpha_{i,\sigma(j)}\) （对 \(1 \leqslant j \leqslant n\)）所得的行列式；若 \(p\) 是在置换 \(\sigma\) 下不变的指标的个数，证明这 \(n\) 个行列式 \(\Delta_i\) ( \(1 \leqslant i \leqslant n\) ）之和等于 \(p\Delta\) （方法与习题 7 相同）。
\item
  设 A 为 n 阶方阵，B 为 A 的一个具有 p 行 q 列的子矩阵，C 为把 A 中属于 B 的每个元素乘以同一标量 \(\alpha\) 后所得的矩阵。证明 det C 的总展开式中的每一项都等于 det A 的总展开式中对应项乘以一个形如 \(\alpha^{r}\) 的标量，其中 \(r \geqslant p + q - n\) ，而 r 取决于所考虑的项（对 det C 作适当的拉普拉斯展开）。
\item
  设 \(\Gamma, \Delta\) 是两个矩阵 \(U\) , \(V\) （\(n\) 阶）的行列式；设 H 和 K 是 \([1, n]\) 的任意两个具有 \(p\) 个元素的子集，而 \((i_k)\) （相应地 \((j_k)\) ）是将 H（相应地，K）的元素按递增顺序排列所得的序列；令 \(\Gamma_{H,K}\) 为在 \(U\) 中把指标为 \(i_k\) 的列换成 \(V\) 的指标为 \(j_k\) 的列（对 \(1 \leqslant k \leqslant p\)）后所得的行列式；类似地，令 \(\Delta_{K,H}\) 为在 \(V\) 中把指标为 \(j_k\) 的列换成 \(U\) 的指标为 \(i_k\) 的列（对 \(1 \leqslant k \leqslant p\)）后所得的行列式。证明：对每个子集 \(H\) （属于 \([1, n]\) ，具有 \(p\) 个元素），
\end{enumerate}

\[
\Gamma \Delta = \sum_ {\mathbf {K}} \Gamma_ {\mathbf {H}, \mathbf {K}} \Delta_ {\mathbf {K}, \mathbf {H}}
\]

其中 K 遍历 \([1, n]\) 的具有 p 个元素的子集（利用第 6 号中的公式 (21) 和 (22)）。

\begin{enumerate}
\def\labelenumi{\arabic{enumi}.}
\setcounter{enumi}{10}
\tightlist
\item
  设 \(\Delta\) 为方阵 \(X\) （\(n\) 阶）的行列式， \(\Delta_p\) 为方阵 \(\bigwedge^p X\) （其阶为 \(\binom{n}{p}\)）的行列式。证明
\end{enumerate}

\[
\Delta_ {p} \Delta_ {n - p} = \Delta_ {(p)} ^ {(n)}
\]

（利用第 6 号中的公式 (21) 和 (22)）。

\begin{enumerate}
\def\labelenumi{\arabic{enumi}.}
\setcounter{enumi}{11}
\tightlist
\item
  矩阵 \((\alpha_{ij})\) （\(n\) 阶）称为中心对称（相应地，斜中心对称）的，如果 \(\alpha_{n - i + 1, n - j + 1} = \alpha_{ij}\) （相应地，\(\alpha_{n - i + 1, n - j + 1} = -\alpha_{ij}\) ）对 \(1 \leqslant i \leqslant n\) , \(1 \leqslant j \leqslant n\) 成立。
\end{enumerate}

\begin{enumerate}
\def\labelenumi{(\alph{enumi})}
\item
  证明 2p 偶数阶中心对称矩阵的行列式可以表示为两个 p 阶行列式的乘积之形式，而 \(2p + 1\) 阶中心对称矩阵的行列式可以表示为一个 p 阶行列式与一个 \(p + 1\) 阶行列式的乘积之形式。
\item
  证明 \(2p\) 偶数阶斜中心对称矩阵的行列式可以表示为两个 \(p\) 阶行列式的乘积之形式。\(2p + 1\) 奇数阶斜中心对称矩阵的行列式，若在 A 中关系 \(2\xi = 0\) 蕴含 \(\xi = 0\) ，则为零；否则它可以表示为 \(\alpha_{p+1,p+1}\) 与两个 \(p\) 阶行列式的乘积之形式。
\end{enumerate}

\begin{enumerate}
\def\labelenumi{\arabic{enumi}.}
\setcounter{enumi}{12}
\tightlist
\item
  设 \(\Delta = \det(\alpha_{ij})\) 是一个 \(n\) 阶行列式， \(\Delta_{ij}\) 是 \(n - 1\) 阶子式，即 \(\alpha_{ij}\) 的余子式。证明
\end{enumerate}

\[
\left| \begin{array}{c c c c c} \alpha_ {1 1} & \alpha_ {1 2} & \ldots & \alpha_ {1 n} & x _ {1} \\ \alpha_ {2 1} & \alpha_ {2 2} & \ldots & \alpha_ {2 n} & x _ {2} \\ \cdot & \cdot & \cdot & \cdot & \cdot \\ \alpha_ {n 1} & \alpha_ {n 2} & \ldots & \alpha_ {n n} & x _ {n} \\ y _ {1} & y _ {2} & \ldots & y _ {n} & z \end{array} \right| = \Delta z - \sum_ {i, j} (- 1) ^ {i + j} \Delta_ {i j} x _ {i} y _ {j}.
\]

若 \(\Delta = 0\) 且诸 \(\alpha_{ij}\) 属于一个域，证明上述行列式是 \(x_1, x_2, \ldots, x_n\) 的一个线性形式与 \(y_1, y_2, \ldots, y_n\) 的一个线性形式的乘积（利用 §5 习题 5 与 II §6 习题 8）。

给出一个例子，说明当标量环 A 不是域时上述结果不再成立（取 A 为环 \(\mathbf{Z} / (6)\) 且 \(n = 2\)）。

\begin{enumerate}
\def\labelenumi{\arabic{enumi}.}
\setcounter{enumi}{13}
\tightlist
\item
  证明恒等式
\end{enumerate}

\[
\left| \begin{array}{cccccc} 0 & 1 & 1 & 1 & \ldots & 1 \\ 1 & 0 & a _ {1} + a _ {2} & a _ {1} + a _ {3} & \ldots & a _ {1} + a _ {n} \\ 1 & a _ {2} + a _ {1} & 0 & a _ {2} + a _ {3} & \ldots & a _ {2} + a _ {n} \\ 1 & a _ {3} + a _ {1} & a _ {3} + a _ {2} & 0 & \ldots & a _ {3} + a _ {n} \\ \cdot & \cdot & \cdot & \cdot & \cdot & \cdot \\ 1 & a _ {n} + a _ {1} & a _ {n} + a _ {2} & a _ {n} + a _ {3} & \ldots & 0 \\ \end{array} \right| = (- 1) ^ {n} 2 ^ {n - 1} \sum_ {i = 1} ^ {n} a _ {1} a _ {2} \ldots a _ {i - 1} a _ {i + 1} \ldots a _ {n}
\]

（利用习题 13）。

\begin{enumerate}
\def\labelenumi{\arabic{enumi}.}
\setcounter{enumi}{14}
\item
  证明：若 A 上的 n 阶方阵的列线性无关，则该矩阵的行也线性无关（参见 II，§10，习题 3）。
\item
  设 E, F 为两个自由 A-模，维数分别为 m 和 n。为使线性映射 \(u: E \rightarrow F\) 是单射，必须且只需 \(m \leqslant n\) ，并且，若 X 表示 u 关于 E 和 F 的任意两组基的矩阵，则不存在标量 \(\mu \neq 0\) ，使得 X 的所有 m 阶子式乘以 \(\mu\) 后均为零。
\end{enumerate}

\begin{enumerate}
\def\labelenumi{\arabic{enumi}.}
\setcounter{enumi}{16}
\tightlist
\item
  设
\end{enumerate}

\[
\sum_ {j = 1} ^ {n} \alpha_ {i j} \xi_ {j} = \beta_ {i} \quad (1 \leqslant i \leqslant m)\tag{*}
\]

是交换环 A 上含 n 个未知数的 m 个线性方程组成的方程组。令 \(x_{j}\) ( \(1 \leqslant j \leqslant n\) ）为该方程组的矩阵 \(X = (\alpha_{ij})\) 的列，而 \(y = (\beta_{i})\) ；假设在 X 中所有阶数 \textgreater p 的子式均为零，但 \(x_{1} \wedge x_{2} \wedge \cdots \wedge x_{p} \neq 0\) 。为使方程组 (*) 有解，必须 \(x_{1} \wedge x_{2} \wedge \cdots \wedge x_{p} \wedge y = 0\) 。反之，若此条件成立，则存在 \(n + 1\) 个属于 A 的元素 \(\xi_{j}\) ( \(1 \leqslant j \leqslant n + 1\) )，使得 \(\xi_{n+1} \neq 0\) 且

\[
\sum_ {j = 1} ^ {n} \alpha_ {i j} \xi_ {j} = \beta_ {i} \xi_ {n + 1} \quad (1 \leqslant i \leqslant m).
\]

\begin{enumerate}
\def\labelenumi{\arabic{enumi}.}
\setcounter{enumi}{17}
\tightlist
\item
  设 E 和 F 分别为维数 m 和 n 的两个自由 A-模， \(u:E \rightarrow F\) 是一个线性映射，而 X 是 u 关于 E 和 F 的任意两组基的矩阵。
\end{enumerate}

\begin{enumerate}
\def\labelenumi{(\alph{enumi})}
\item
  若 \(m \geqslant n\) 并且存在 \(X\) 的一个 \(n\) 阶可逆子式，则 \(u\) 是满射。
\item
  反之，若 \(u\) 是满射，证明 \(m \geqslant n\) ，且存在 \(X\) 的一个非零的 \(n\) 阶子式。此外，若在环 \(A\) 中，由不可逆元素生成的理想异于 \(A\) ，则存在 \(X\) 的一个 \(n\) 阶可逆子式（考虑外幂 \(\bigwedge^{n}u\) ）。
\item
  设 B 是一个交换环，A 是乘积环 \(B \times B\) 。给出一个从 A-模 \(A^{2}\) 到 A 上的满射线性映射的例子，其矩阵（关于 \(A^{2}\) 与 A 的典范基）的所有元素都是零因子。
\end{enumerate}

\begin{enumerate}
\def\labelenumi{\arabic{enumi}.}
\setcounter{enumi}{18}
\item
  设 X 是交换域上的一个矩阵。要使 X 的秩为 p，只需存在 X 的一个 p 阶子式，它 \(\neq0\) ，并且所有包含这个 p 阶子式的 \(p+1\) 阶子式均为零（证明此时 X 的每一列都是所述 p 阶子式的元素所属的那 p 列的线性组合）。
\item
  设 \(\mathbf{A}\) 为一个交换环， \(\mathbf{M}\) 为任意一个 A-模，而 \(U = (\alpha_{ij})\) 为一个类型为 \((m, n)\) 、元素属于 \(\mathbf{A}\) 的矩阵。
\end{enumerate}

\begin{enumerate}
\def\labelenumi{(\alph{enumi})}
\item
  假设 \(m = n\) ，并设 \(\Delta = \det(U)\) ；那么，若 \(x_1, \ldots, x_n\) 是 \(M\) 的元素，使得 \(\sum_{j=1}^{n} \alpha_{ij} x_j = 0\) （对 \(1 \leqslant i \leqslant n\) ），则 \(\Delta x_i = 0\) 对每个指标 \(i\) 成立。
\item
  假设 \(m\) 和 \(n\) 是任意的， \(U\) 的所有阶 \(> r\) 的子式均为零，且存在一个 \(r\) 阶子式是可逆的，例如矩阵 \((\alpha_{ij})\) （其中 \(i \leqslant r\) 且 \(j \leqslant r\) ）的子式；那么，若 \(a_i\) 表示指标为 \(i\) 的行（在 \(U\) 中），则这些 \(a_i\) 中指标 \(\leqslant r\) 的那些构成 \(A^n\) 的由 \(U\) 的所有行生成的子模的一组基；因而对于 \(i > r\) ，有 \(a_i = \sum_{k=1}^{r} \lambda_{ik} a_k\) 。然后设 \(y_1, \ldots, y_m\) 是 \(M\) 的元素；为使存在元素 \(x_1, \ldots, x_n\) （属于 \(M\) ）满足方程组 \(\sum_{j=1}^{n} \alpha_{ij} x_j = y_i\) （对 \(1 \leqslant i \leqslant m\) ），必要且充分的条件是 \(y_i = \sum_{k=1}^{r} \lambda_{ik} y_k\) （对 \(i > r\) ）。
\end{enumerate}

¶ 21. 设 K 为一个无限交换域，m、n、r 为三个整数 \(\geqslant0\) ，使得 \(r\leqslant n\leqslant m\) ；设 \(\mathbf{M}_{m,n}(\mathbf{K})\) 为由 K 上类型为 \((m,n)\) 的矩阵组成的（mn 维）K-向量空间，而 V 为 \(\mathbf{M}_{m,n}(\mathbf{K})\) 的一个向量子空间，使得 r 是矩阵 \(X\in V\) 的秩的最大值。我们打算证明不等式

\[
\dim (\mathbf {V}) \leqslant m r.\tag{*}
\]

\begin{enumerate}
\def\labelenumi{(\alph{enumi})}
\tightlist
\item
  证明我们可以假设 \(m = n\) 且矩阵
\end{enumerate}

\[
X _ {0} = \left( \begin{array}{c c} I _ {r} & 0 \\ 0 & 0 \end{array} \right)
\]

属于 V。由此推出每个矩阵 \(X \in \mathbf{V}\) 都具有形式

\[
\left( \begin{array}{c c} X _ {1 1} & X _ {1 2} \\ X _ {2 1} & 0 \end{array} \right)
\]

（其中 \(X_{11}\) 的阶为 \(r\) ）且满足条件 \(X_{21}X_{12} = 0\) （利用如下事实：对所有 \(\xi \in \mathbf{K}\) ，矩阵 \(\xi X_0 + X\) 的秩 \(\leqslant r\) ）。

\begin{enumerate}
\def\labelenumi{(\alph{enumi})}
\setcounter{enumi}{1}
\tightlist
\item
  从 (a) 推出，对于两个矩阵 \(X\) , \(Y\) （属于 \(\mathbf{V}\) ），
\end{enumerate}

\[
X _ {2 1} Y _ {1 2} + Y _ {2 1} X _ {1 2} = 0.
\]

\begin{enumerate}
\def\labelenumi{(\alph{enumi})}
\setcounter{enumi}{2}
\tightlist
\item
  对每个矩阵 \(X \in \mathbf{V}\) ，我们记 \(F(X) = (X_{11}X_{12}) \in \mathbf{M}_{r,n}(\mathbf{K})\) ；\(f\) 是一个线性映射，且 \(\dim(\mathbf{V}) = \dim(\operatorname{Im}(f)) + \dim(\operatorname{Ker}(f))\) 。另一方面，设 \(u_x\) 表示线性形式 \((\mathrm{Y}_{11}\mathrm{Y}_{12}) \mapsto \operatorname{Tr}(\mathrm{X}_{21}\mathrm{Y}_{12})\) （定义在 \(\mathbf{M}_{r,n}(\mathbf{K})\) 上）。证明线性映射 \(X \mapsto u_X\) （从 \(\operatorname{Ker}(f)\) 到对偶 \((\mathbf{M}_{r,n}(\mathbf{K}))^*\) 中）是单射；并通过注意在 \(X \mapsto u_X\) 下 \(\operatorname{Ker}(f)\) 的像包含于 \(\operatorname{Im}(f)\) 的正交补中来完成 (*) 的证明。
\end{enumerate}

¶ 22. 设 F 是从 \(\mathbf{M}_{n}(\mathrm{K})\) 到一个集合 E 中的映射，其中 K 是一个交换域，使得对每三个矩阵 \(X\) , \(Y\) , \(Z\) （属于 \(\mathbf{M}_{n}(\mathrm{K})\) ），

\[
\mathbf {F} (X Y Z) = \mathbf {F} (X Z Y).
\]

我们排除 \(\mathbf{M}_{2}(\mathbf{F}_{2})\) 的情形；然后证明存在一个从 K 到 E 中的映射 \(\Phi\) ，使得 \(\mathrm{F}(X)=\Phi(\det(X))\) 。（利用命题 17（第 9 号）的推论，首先证明，对于 \(U\in\mathbf{SL}_{n}(\mathrm{K})\) ，有 \(\mathrm{F}(X)=\mathrm{F}(XU)\) ；由此推出，在 \(\mathbf{GL}_{n}(\mathrm{K})\) 中， \(\mathrm{F}(X)\) 仅依赖于 \(\det(X)\) 。另一方面，若对于 \(r\leqslant n-1\) ，我们记

\[
X _ {r} = \left( \begin{array}{c c} I _ {r} & 0 \\ 0 & 0 \end{array} \right)
\]

，证明存在一个矩阵 \(Y\) （秩为 \(r\) ），使得 \(X_{r-1} = YX_r\) 且 \(Y = X_rY\) ，并由此推出 \(F(X)\) 对所有秩 \(< n\) 的矩阵取相同的值。

\begin{enumerate}
\def\labelenumi{\arabic{enumi}.}
\setcounter{enumi}{22}
\tightlist
\item
  在每个A-代数E中，若 \(x_{1}, \ldots, x_{n}\) 是E的元素，我们记作
\end{enumerate}

\[
\left[ x _ {1} x _ {2} \dots x _ {n} \right] = \sum_ {\sigma \in \mathfrak {S} _ {n}} \varepsilon_ {\sigma} x _ {\sigma (1)} x _ {\sigma (2)} \dots x _ {\sigma (n)}.
\]

证明：若 E 在 A 上具有有限生成系，则存在整数 N，使得 \([x_{1}x_{2}\ldots x_{N}] = 0\) 对 E 的所有元素 \(x_{j} (1 \leqslant j \leqslant N)\) 成立。

\begin{enumerate}
\def\labelenumi{\arabic{enumi}.}
\setcounter{enumi}{23}
\tightlist
\item
  设 \(X_{i}\) ( \(1 \leqslant i \leqslant m\) ) 是交换环 A 上具有相同阶 \(n\) 的方阵。证明若 \(m\) 是偶数，则（习题 23）
\end{enumerate}

\[
\operatorname{Tr} \left[ X _ {1} X _ {2} \dots X _ {m} \right] = 0
\]

且若 \(m\) 是奇数，则

\[
\operatorname{Tr} \left[ X _ {1} X _ {2} \dots X _ {m} \right] = m. \operatorname{Tr} \left(X _ {m} \left[ X _ {1} X _ {2} \dots X _ {m - 1} \right]\right).
\]

（若 \(\lambda\) 是循环置换 \((1,2,\ldots,m)\) ，C 是 \(\mathfrak{S}_m\) 的由 \(\lambda\) 生成的循环子群，而 H 是 \(\mathfrak{S}_m\) 的由保持 \(m\) 不变的置换组成的子群，注意 \(\mathfrak{S}_m\) 中的每个置换都可以唯一地写成 \(\sigma\tau\) ，其中 \(\sigma\in\mathrm{H}\) 且 \(\tau\in\mathrm{C}\) ；然后利用矩阵乘积的迹在因子的循环置换下不变这一事实。）

¶ 25. 设 A 是一个包含有理数域 Q 的交换环，并设 X 是 A 上一个偶数阶 \(2n \geqslant 2\) 的方阵。设 L 为子集 \(\{1, 2, \ldots, n + 1\}\) （属于 \(\{1, 2, \ldots, 2n\}\) ），而 L' 为其补集。证明以下公式（“Redei 恒等式”）：

\[
(*) \quad \det (X _ {\mathrm{L}, \mathrm{L}}) \det (X _ {\mathrm{L} ^ {\prime}, \mathrm{L} ^ {\prime}}) = \sum_ {\mathrm{H}, \mathrm{K}} (- 1) \frac {r - 1}{n \binom {n - 1} {r}} \det (X _ {\mathrm{H}, \mathrm{K}}) \det (X _ {\mathrm{H} ^ {\prime}, \mathrm{K} ^ {\prime}})
\]

其中 H 遍历 \(\{1,2,\ldots,2n\}\) 的包含 l 的 n 元子集组成的集合，而 K 遍历 L 的 n 元子集组成的集合，且 \(r=\operatorname{Card}(H\cap L')\) 。（计算 (*) 右端中项 \(x_{\sigma(1),1}x_{\sigma(2),2}\ldots x_{\sigma(2n),2n}\) 对任意置换 \(\sigma\in S_{2n}\) 的系数，注意该系数仅可能 \(\neq0\) （当 \(H=\sigma(K)\) ）；证明它仅可能 \(\neq0\) （当 \(\sigma(L)=L\) ）。）

\begin{enumerate}
\def\labelenumi{\arabic{enumi}.}
\setcounter{enumi}{25}
\tightlist
\item
  \begin{enumerate}
  \def\labelenumii{(\alph{enumii})}
  \tightlist
  \item
    在命题 19（第 10 号）的记号下，假设存在两个 A{[}X{]}-线性映射 \(g_1, g_2\) （从 M{[}X{]} 到 M'{[}X{]} 中），使得 \(\psi' \circ g_2 = g_1 \circ \psi\) ；证明此时存在一个 A{[}X{]}-线性映射 \(g\) （从 \(M_u\) 到 \(M_{u'}\) 中），使得 \(\phi' \circ g_1 = \phi' \circ \bar{g}\) 。进一步，若 \(g_1\) 和 \(g_2\) 是 A{[}X{]}-同构（从 M{[}X{]} 到 M'{[}X{]} 上），则 \(g\) 是一个 A{[}X{]}-同构（从 \(M_u\) 到 \(M_{u'}\) 上）。
  \end{enumerate}
\end{enumerate}

\begin{enumerate}
\def\labelenumi{(\alph{enumi})}
\setcounter{enumi}{1}
\tightlist
\item
  在第 10 号的记号下，自同态 \(u, u'\) 称为等价的，如果存在两个从 M 到 M 上的同构 \(f_1, f_2\) ，使得 \(u' \circ f_2 = f_1 \circ u\) ；若存在一个从 M 到 M' 上的同构 \(f\) ，使得 \(u' \circ f = f \circ u\) ，则称它们为相似的。由 (a) 推出，要使 \(u\) 和 \(u'\) 相似，当且仅当自同态 \(X - \bar{u}\) 和 \(X - \bar{u}'\) （分别为 A{[}X{]}-模 M{[}X{]} 和 M'{[}X{]} 的自同态）等价。
\end{enumerate}

\BourbakiExercisesSection

\begin{enumerate}
\def\labelenumi{\arabic{enumi}.}
\tightlist
\item
  设 K 是至少含有 3 个元素的交换域，A 是 K 上的一个代数，其基由单位元 1 和两个元素 \(e_1, e_2\) 组成，使得 \(e_1^2 = e_1, e_1e_2 = e_2, e_2e_1 = e_2^2 = 0\) ；令 \(\mathbf{A}^0\) 为 A 的反代数。证明在 A 中存在元素 \(x\) ，使得 \(\operatorname{Tr}_{\mathbf{A} / \mathbf{K}}(x) \neq \operatorname{Tr}_{\mathbf{A}^0 / \mathbf{K}}(x)\) 且 \(\mathbf{N}_{\mathbf{A} / \mathbf{K}}(x) \neq \mathbf{N}_{\mathbf{A}^0 / \mathbf{K}}(x)\) 。
\end{enumerate}

\BourbakiExercisesSection

\begin{enumerate}
\def\labelenumi{\arabic{enumi}.}
\tightlist
\item
  若 A 和 B 是 K-代数， \(\alpha\) 和 \(\beta\) 是从 A 到 B 中的两个 K-同态，则 A 到 B 中的一个 \((\alpha, \beta)\) -导子是指满足关系式
\end{enumerate}

\[
d (x y) = (d x) \alpha (y) + \beta (x) (d y) \quad \text { for } \quad x, y \quad \text { in } \quad A,
\]

换言之，即定义 1（第 2 号）意义下的导子，其中 \(\mathbf{A} = \mathbf{A}' = \mathbf{A}''\) , \(\mathbf{B} = \mathbf{B}' = \mathbf{B}''\) , \(d = d' = d''\) , \(\mu: \mathbf{A} \times \mathbf{A} \to \mathbf{A}\) 是乘法， \(\lambda_1: \mathbf{B} \times \mathbf{A} \to \mathbf{B}\) 是 K-双线性映射 \((z, x) \mapsto z\alpha(x)\) ，而 \(\lambda_2: \mathbf{A} \times \mathbf{B} \to \mathbf{B}\) 是 K-双线性映射 \((x, z) \mapsto \beta(x)z\) 。

从此以后，设 A 和 B 为交换域，且 \(\alpha \neq \beta\) 。证明此时存在一个元素 \(b \in \mathbf{B}\) ，使得 \(d\) 是 \((\alpha, \beta)\) -导子 \(x \mapsto b\alpha(x) - \beta(x)b\) 。（可以假设 \(d \neq 0\) 。首先证明 \(d\) 的核 N 是由 \(x \in \mathbf{A}\) （满足 \(\alpha(x) = \beta(x)\) ）组成的集合，然后证明对于 \(x \notin \mathbf{N}\) ，元素 \(b = d(x)/( \alpha(x) - \beta(x))\) 与 \(x\) 的选取无关。）

\begin{enumerate}
\def\labelenumi{\arabic{enumi}.}
\setcounter{enumi}{1}
\tightlist
\item
  设 K 是一个特征 \(\neq2\) 的交换域，A 和 B 为两个 K-代数，且 \(\alpha\) 为从 A 到 B 中的一个 K-线性映射；A 到 B 中的一个 \(\alpha\) -导子是指从 A 到 B 中的一个 K-线性映射 d，它满足关系式
\end{enumerate}

\[
d (x y) = (d x) \alpha (y) + \alpha (x) (d y) \quad \text { for } \quad x, y \quad \text { in } \quad A,
\]

换言之，即定义 1 意义下的一个导子，其中 \(\mathbf{A} = \mathbf{A}' = \mathbf{A}''\) , \(\mathbf{B} = \mathbf{B}' = \mathbf{B}''\) , \(d = d' = d''\) , \(\mu\) 是 \(\mathbf{A}\) 中的乘法， \(\lambda_1\) 是映射 \((z,x) \mapsto z\alpha(x)\) （从 \(\mathbf{B} \times \mathbf{A}\) 到 \(\mathbf{B}\) 中），而 \(\lambda_2\) 是映射 \((x,z) \mapsto \alpha(x)z\) （从 \(\mathbf{A} \times \mathbf{B}\) 到 \(\mathbf{B}\) 中）。

证明：若 B 是交换域、K 的一个扩域，且 \(d \neq 0\) ，则存在 B 中的 \(b \neq 0\) ，使得

\[
\alpha (x y) = \alpha (x) \alpha (y) + b (d x) (d y) \quad \text { for } \quad x, y \quad \text { in } \quad A.
\]

若 \(b = c^2\) ，其中 \(c \in \mathbf{B}\) , \(c \neq 0\) ，则存在两个同态 \(f, g\) （从 \(\mathbf{A}\) 到 \(\mathbf{B}\) 中），使得 \(\alpha(x) = (f(x) + g(x)) / 2\) , \(dx = (f(x) - g(x)) / 2c\) 。否则，存在一个二次代数（§ 2，第 3 号） \(\mathbf{B}'\) （在 \(\mathbf{B}\) 上），一个 \(\mu \in \mathbf{B}'\) ，使得 \(\mu^2 = c\) ，以及一个同态 \(f\) （从 \(\mathbf{A}\) 到 \(\mathbf{B}'\) 中），使得 \(\alpha(x) = (f(x) + \overline{f(x)}) / 2\) , \(dx = (f(x) - \overline{f(x)}) / 2\) 。

¶ 3. 设 K 为一个域，无论交换与否，E 为左向量空间 \(\mathrm{K}_{s}^{(\mathbf{N})}\) ，并设 \((e_{n})_{n\in\mathbf{N}}\) 为该空间的典范基。另一方面，设 \(\sigma\) 是 K 的一个自同态，且 d 是 K 到自身的一个 \((\sigma,1_{\mathbf{K}})\) -导子（习题 1），换言之，即 K 到自身的一个加法映射，使得

\[
d (\xi \eta) = (d \xi) \sigma (\eta) + \xi (d \eta).
\]

E 的元素是形如 \(\alpha_{n}e_{n}\) 的乘积的线性组合，并且在 E 上定义了一个乘法（E × E 到 E 中的 Z-双线性映射），其条件为

\[
\begin{array}{l} (\alpha e _ {n}) (\beta e _ {m}) = (\alpha e _ {n - 1}) (\sigma (\beta) e _ {m + 1} + (d \beta) e _ {m}) \text { for } n \geqslant 1 \\ (\alpha e _ {0}) (\beta e _ {m}) = (\alpha \beta) e _ {m}. \end{array}
\]

因此，在 E 上定义了一个（一般是非交换的）环结构，以 \(e_{0}\) 为单位元（与 K 的元素 l 等同）。若记 \(X = e_{1}\) ，则 \(e_{n} = X^{n}\) 对于 \(n \geqslant 1\) ，从而 E 的元素可以写成

\[
\alpha_ {0} + \alpha_ {1} X + \dots + \alpha_ {n} X ^ {n},
\]

根据乘法规则

\[
\mathrm{X} \alpha = \sigma (\alpha) \mathrm{X} + (d \alpha) \quad \text { for } \quad \alpha \in \mathrm{K}.
\]

这个环记作 K{[}X; σ, d{]} ，并称为关于 σ 和 d 的、系数在 K 中的 X 的非交换多项式环。~对于 E 中的一个非零元素 \(z = \sum_{j} \alpha_{j} X^{j}\) ，次数 deg(z) 是使得 \(\alpha_{j} \neq 0\) 的最大整数 j。

\begin{enumerate}
\def\labelenumi{(\alph{enumi})}
\tightlist
\item
  证明：如果 \(u, v\) 是两个元素 \(\neq 0\) （属于 \(E\) ），那么：
\end{enumerate}

\[
\begin{array}{r l} u v \neq 0 & \text { and } \quad \deg (u v) = \deg (u) + \deg (v); \\ \deg (u + v) & \leqslant \sup (\deg (u), \deg (v)) \quad \text { if } u + v \neq 0. \end{array}
\]

此外，若 \(\deg(u) \geqslant \deg(v)\) ，则存在 E 中的两个元素 w, r ，使得 \(u = wv + r\) ，并且要么 r = 0，要么 \(\deg(r) < \deg(v)\) （E 中的“欧几里得除法”）。

\begin{enumerate}
\def\labelenumi{(\alph{enumi})}
\setcounter{enumi}{1}
\tightlist
\item
  反之，设 R 是一个环，且存在一个映射 \(u \mapsto \deg(u)\) ，它把集合 \(R'\) （由 R 中 \(\neq 0\) 的元素组成）映入 N，并满足 (a) 的条件。证明由 0 与非零元素 \(u \in R\) （满足 \(\deg(u) = 0\) ）组成的集合 K 是一个（未必交换的）域。若 \(K \neq R\) 且 \(x\) 是 R 中的一个元素，使得整数 \(\deg(x) > 0\) 为尽可能小的，证明
\end{enumerate}

R 的每个元素都可以唯一地写成形式 \(\sum_{i} \alpha_j x^j\) ，其中 \(\alpha_j \in K\) 。（要看出 R 的每个元素都具有这种形式，可用反证法论证，考虑一个 \(u \in R\) （不具有这种形式的元素），且对它而言 \(\deg(u)\) 是尽可能小的。）证明存在 K 的一个自同态 \(\sigma\) 和 K 的一个 \((\sigma, 1_K)\) -导子 \(d\) ，使得 \(x\alpha = \sigma(\alpha)x + d\alpha\) 对所有 \(\alpha \in K\) 成立。由此推出 R 同构于一个域或一个环 K{[}X; \(\sigma, d]\) 。

\begin{enumerate}
\def\labelenumi{\arabic{enumi}.}
\setcounter{enumi}{3}
\tightlist
\item
  设 M 是类型为 N 的分次 K-模；证明每个次数为 r 的分次自同态都可以唯一地延拓为泛反交换代数 \(\mathbf{G}(\mathbf{M})\) （§ 7，习题 13）的次数为 r 的一个导子。
\end{enumerate}

* 5. 设 K 是特征为 \(p > 0\) 的域，B 是单不定元的有理函数域 K(X)，而 A 是 B 的子域 K(X \(^{p}\) ) 。证明典范同态 \(\Omega_{\mathbf{K}}(\mathbf{A}) \otimes_{\mathbf{A}} \mathbf{B} \to \Omega_{\mathbf{K}}(\mathbf{B})\) 不是单射。*

\BourbakiExercisesSection

\begin{enumerate}
\def\labelenumi{\arabic{enumi}.}
\item
  若 V 是有限生成投射 A-模，则 End(V) 亦然，且典范地等同于 \(V \otimes_{A} V^{*}\) （II，§4，第 2 号，命题 2 的推论）。若将每个元素 \(u \in \text{End}(V)\) 与 End(V) 上的 A-线性形式 \(v \mapsto \text{Tr}(uv)\) 相关联，就定义了从 End(V) 到其对偶 (End(V))* 上的一个 A-线性双射。在此映射下将 End(V) 与其对偶 A-模等同，则 End(V) 上的 A-代数结构通过对偶性（第 1 号，例 3）在 End(V) 上定义一个 A-余代数结构，它是余结合的，并以映为 \(V \to A\) 的形式 Tr 作为余单位。特别地，对 \(V = A^{n}\) ，在 \(\mathbf{M}_{n}(\mathbf{A})\) 上定义了一个余代数结构，其余积由 \(c(E_{ij}) = \sum_{k} E_{kj} \otimes E_{ik}\) 给出。这个余代数结构与 \(\mathbf{M}_{n}(\mathbf{A})\) 上的代数结构不定义双代数结构。
\item
  证明在定义 3（第 4 号）中，条件（3）（关于分次双代数的）等价于：乘积 \(m: \mathbf{E} \otimes \mathbf{E} \to \mathbf{E}\) 是从分次余代数 \(\mathbf{E} \otimes \mathbf{E}\) 到分次余代数 \(\mathbf{E}\) 中的态射。
\item
  证明：若 \(\mathbf{M}\) 是一个 A-模，则在 \(\mathsf{T}(\mathbf{M})\) 上存在且仅存在一个余交换双代数结构，其代数结构是通常的结构，且其余积 \(c\) 满足：对所有 \(x\in \mathbf{M}\) ， \(c(x) = 1\otimes x + x\otimes 1\) 。
\item
  设 \(\mathbf{E}\) 是 \(\mathbf{A}\) 上的一个交换双代数， \(m\) 是乘积 \(\mathbf{E} \otimes \mathbf{E} \to \mathbf{E}\) ， \(c\) 是余积 \(\mathbf{E} \to \mathbf{E} \otimes \mathbf{E}\) ， \(e\) 是单位元，而 \(\gamma\) 是 \(\mathbf{E}\) 的余单位。
\end{enumerate}

\begin{enumerate}
\def\labelenumi{(\alph{enumi})}
\tightlist
\item
  证明：对于每一个交换A-代数B，该合成法则将每一个有序对
\end{enumerate}

\[
(u, v) \in \operatorname{Hom} _ {\mathbf {A} - \mathrm{alg.}} (\mathbf {E}, \mathbf {B}) \times \operatorname{Hom} _ {\mathbf {A} - \mathrm{alg.}} (\mathbf {E}, \mathbf {B})
\]

对应到 A-代数同态

\[
\mathrm{E} \xrightarrow {c} \mathrm{E} \otimes \mathrm{E} \xrightarrow {u \otimes v} \mathrm{B} \otimes \mathrm{B} \xrightarrow {m _ {\mathrm{B}}} \mathrm{B}
\]

（其中 \(m_{\mathbf{B}}\) 是 B 中的乘积）在 \(\operatorname{Hom}_{\mathbf{A}\text{-alg.}}(\mathbf{E}, \mathbf{B})\) 上定义了一个幺半群结构。（b）为使对每个交换 A-代数 B， \(\operatorname{Hom}_{\mathbf{A}\text{-alg.}}(\mathbf{E}, \mathbf{B})\) 都是一个群，必要且充分的是存在一个 A-代数同态 \(i: \mathbf{E} \to \mathbf{E}\) ，使得 \(i(e) = e\) ，且复合映射 \(m \circ (1_{\mathbf{E}} \otimes i) \circ c\) 与 \(m \circ (i \otimes 1_{\mathbf{E}}) \circ c\) 都等于映射 \(x \mapsto \gamma(x)e\) 。这时称 \(i\) 为 E 中的一个逆；它是从 E 到相反双代数 \(E^0\) 上的一个同构。

\begin{enumerate}
\def\labelenumi{\arabic{enumi}.}
\setcounter{enumi}{4}
\item
  设 \(E_{1}\) , \(E_{2}\) 是 A 上的两个双代数。证明在张量积 \(E_{1} \otimes_{A} E_{2}\) 上， \(E_{1}\) 和 \(E_{2}\) 上的代数与余代数结构的张量积在 \(E_{1} \otimes_{A} E_{2}\) 上定义了一个双代数结构。
\item
  设 M 是类型为 N 的分次 A-模。在泛反交换代数 G(M) 上定义一个斜分次双代数结构（§ 7，习题
\end{enumerate}

13）；由此导出一个典范的分次代数同态 \(\mathbf{G}(\mathbf{M}^{*}) \to \mathbf{G}(\mathbf{M})^{*\mathbf{gr}}\) ，以及 \(\mathbf{G}(\mathbf{M})\) 的元素与 \(\mathbf{G}(\mathbf{M}^{*})\) 的元素之间的内积概念。

\begin{enumerate}
\def\labelenumi{\arabic{enumi}.}
\setcounter{enumi}{6}
\tightlist
\item
  在命题 11（第 9 号）的假设和记号下，证明公式
\end{enumerate}

\[
(\det g) \phi^ {- 1} = \bigwedge (g) \circ \phi^ {- 1} \circ \bigwedge (^ {t} g).
\]

由该公式与 (79)（第 11 号）可导出用余子式矩阵表示方阵的逆的表达式的一个新证明（§ 8，第 6 号，公式 (26)）。

\begin{enumerate}
\def\labelenumi{\arabic{enumi}.}
\setcounter{enumi}{7}
\tightlist
\item
  设 E 是维数为 n 的自由 A-模。~对从 \(\bigwedge(\mathbf{M})\) 到一个 A-模 G 中的每个 A-线性映射 v，我们可以写成 \(v = \sum_{p=0}^{n} v_{p}\) ，其中 \(v_{p}\) 是 v 在 \(\bigwedge^{p}(\mathbf{M})\) 上的限制；然后我们设 \(\eta v = \sum_{p=0}^{n} (-1)^{p} v_{p}\) ，使得 \(\eta^{2} v = v\) 。设 u 是从 M 到其对偶 \(M^{*}\) 上的一个同构，并设 \(\Delta\) 是 u 关于 M 的一组基 \((e_{i})\) 及 \(M^{*}\) 的对偶基的矩阵的行列式（该行列式不依赖于 M 的基的选取）。若 \(\phi\) 是从 \(\bigwedge(\mathbf{M})\) 到 \(\bigwedge(\mathbf{M}^{*})\) 上的、由 \(e = e_{1} \wedge e_{2} \wedge \cdots \wedge e_{n}\) 出发定义的同构，证明公式
\end{enumerate}

\[
\eta^ {n + 1} \Delta \phi = \bigwedge (t u) \circ \phi^ {- 1} \circ \bigwedge (u).
\]

\begin{enumerate}
\def\labelenumi{\arabic{enumi}.}
\setcounter{enumi}{8}
\tightlist
\item
  方阵 \(X\) （\(n\) 阶，在 \(A\) 上）的伴随矩阵是矩阵 \(\tilde{X} = (\det(X^{it}))\) （由 \(X\) 的 \(n - 1\) 阶子式组成）。证明若 \(X\) 可逆，则 \(\det(\tilde{X}) = (\det X)^{n-1}\) ，且每个子式 \(\det(\tilde{X}_{H,K})\) （\(\tilde{X}\) 的 \(p\) 阶子式）由公式给出
\end{enumerate}

\[
\det (\tilde {X} _ {\mathrm{H}, \mathrm{K}}) = (\det X) ^ {p - 1} (X _ {\mathrm{H} ^ {\prime}, \mathrm{K} ^ {\prime}})\tag{1}
\]

其中 H' 和 K' 分别是 H 和 K 在 \([1, n]\) 中的补集（“雅可比恒等式”；利用第 11 号的关系式 (80)）。

\begin{enumerate}
\def\labelenumi{\arabic{enumi}.}
\setcounter{enumi}{9}
\tightlist
\item
  从每个恒等式 \(\Phi = 0\) （任意可逆方阵 \(X\) （\(n\) 阶，在 A 上）的子式之间的）可以导出另一个恒等式 \(\tilde{\Phi} = 0\) ，称为 \(\Phi = 0\) 的余恒等式，其做法是将恒等式 \(\Phi = 0\) 应用于伴随矩阵 \(\tilde{X}\) （\(X\) 的）的子式，然后利用习题 9 中的恒等式 (1) 把 \(\tilde{X}\) 的子式替换为 \(X\) 的子式的函数。用这种方法证明以下恒等式：
\end{enumerate}

\[
\det (X ^ {i h}) \det (X ^ {j k}) - \det (X ^ {i k}) \det (X ^ {j h}) = (\det X) \det (X ^ {i j, h k}) \quad \text { for } \quad i <   j
\]

其中 \(X^{ij, hk}\) 表示 X 的 n-2 阶子式，它是通过在 X 中删去指标为 i、j 的行以及指标为 h、k 的列而得到的。

¶ 11. 设 \(\Phi = 0\) 是交换环 A 上一个可逆方阵 \(X\) （\(n\) 阶）的子式之间的恒等式， \(\tilde{\Phi} = 0\) 为其余恒等式（习题 10），而 \(k\) 为一个整数 \(>0\) 。设 \(Y\) 是 \(n + k\) 阶可逆方阵，并设 \(\tilde{Y}_0\) 是伴随矩阵 \(\tilde{Y}\) 通过在 \(\tilde{Y}\) 中删去指标 \(\leqslant k\) 的行和指标 \(\leqslant k\) 的列所成的子矩阵。若假设 \(\tilde{Y}_0\) 可逆，并将恒等式 \(\tilde{\Phi} = 0\) 应用于 \(\tilde{Y}_0\) 的子式，然后利用习题 9 的恒等式 (1)，把该恒等式中出现的每个子式（视为 \(\tilde{Y}\) 的子式）替换为用 \(Y\) 的子式表示的表达式，就得到一个恒等式 \(\Phi_k = 0\) （在矩阵 \(Y\) 的子式之间，当 \(Y\) 和 \(\tilde{Y}_0\) 可逆时成立），称为 \(k\) 阶的扩张，即恒等式 \(\Phi = 0\) 的扩张。

特别地，设 \(A = (\alpha_{ij})\) 是一个 \(n + k\) 阶可逆方阵，B 是 A 的 k 阶子矩阵（在 A 中删去指标 \textgreater k 的行和列而得），而 \(\Delta_{ij}\) 是 \(k + 1\) 阶矩阵的行列式（在 A 中删去指标 \textgreater k 但保留指标 \(k + i\) 的行、以及指标 \textgreater k 但保留指标 \(k + j\) 的列而得）；若 C 表示 n 阶矩阵 \((\Delta_{ij})\) ，证明恒等式

\[
\det C = (\det A) (\det B) ^ {n - 1}
\]

当 \(A\) 和 \(B\) 可逆时成立（证明它是 \(k\) 阶的扩张，施于 \(n\) 阶行列式的完全展开）。

叙述通过扩张拉普拉斯展开（§ 8，第 6 号）以及 § 8 习题 2 的恒等式所得到的恒等式。

¶ 12. 设 \(X = (\xi_{ij})\) 是一个 \(n\) 阶方阵（在交换域 \(K\) 上）；令 \(H\) 是 \([1, n]\) 的具有 \(p\) 个元素的一个子集，而 \(H'\) 为 \(H\) 在 \([1, n]\) 中的补集；假设对每对有序指标 \(h \in H\) , \(k \in H'\) ，

\[
\sum_ {i} \xi_ {i h} \xi_ {i k} = 0.
\]

证明：对 \([1, n]\) 中具有 p 个元素的所有子集 L, M ，有

\[
\rho_ {\mathbf {M}, \mathbf {M} ^ {\prime}} \det (X _ {\mathrm{L}, \mathrm{H}}) \det (X _ {\mathbf {M} ^ {\prime}, \mathbf {H} ^ {\prime}}) - \rho_ {\mathrm{L}, \mathrm{L} ^ {\prime}} \det (X _ {\mathbf {M}, \mathbf {H}}) \det (X _ {\mathrm{L}, \mathrm{H} ^ {\prime}}) = 0.
\]

（把列 \(x_h\) （指标 \(h \in \mathbf{H}\) ，在 \(X\) 中）看作 \(\mathbf{E} = \mathbf{K}^n\) 中的向量，把列 \(x_k^*\) （指标 \(k \in \mathbf{H}'\) ）看作 \(\mathbf{E}^*\) 中的向量，并证明 \((n - p)\) -形式 \(x_{\mathbf{H}'}^*\) 与 \(\phi_p(x_{\mathbf{H}})\) 成比例。）

\begin{enumerate}
\def\labelenumi{\arabic{enumi}.}
\setcounter{enumi}{12}
\tightlist
\item
  设 \(\Gamma\) 和 \(\Delta\) 是交换环上两个 n 阶可逆矩阵的行列式，且 \(\Gamma_{ij}\) 是通过将 \(\Gamma\) 中第 i 列替换为 \(\Delta\) 的第 j 列所得到的行列式。证明
\end{enumerate}

\[
\det (\Gamma_ {i j}) = \Gamma^ {n - 1} \Delta .
\]

（把 \(\Gamma_{ij}\) 按第 \(i\) 列展开，并利用习题 9。）

\begin{enumerate}
\def\labelenumi{\arabic{enumi}.}
\setcounter{enumi}{13}
\item
  设 M 为维数 n \textgreater{} 1 的自由 A-模。证明：每个同构 \(\psi\) （从 \(\bigwedge^{p}(M)\) 到 \(\bigwedge^{n-p}(M^{*})\) 上），使得 \(\bigwedge^{n-p}(\check{u}) \circ \psi = \psi \circ \bigwedge^{p}(u)\) 对 M 的每个行列式等于 1 的自同构 u 成立，都是第 11 号命题 12 中定义的同构 \(\phi_{p}\) 之一（按 § 7 习题 4 的方法进行）。
\item
  设 E 为维数为 n 的自由 A-模。~设 z 为 E 上的一个 p-向量， \(z^{*}\) 为 E 上的一个 \((p + q)\) -形式；z 可按典范方式（§ 7 习题 8）等同于 \(az_{1}\) （一个 p 阶反变张量 \(z_{1}\) 的斜对称化），而 \(z^{*}\) 等同于一个 \(p + q\) 阶协变张量的斜对称化。如果在混合张量 \(z_{1}\) . \(z^{*}\) 中，第 k 个反变指标与第 \((p + k)\) 个协变指标对 \(1 \leqslant k \leqslant p\) 进行缩并，证明如此得到的 q 阶协变张量是斜对称化的，并且除一个符号外典范地等同于 q-形式 \(z^{*} \perp z\) 。
\item
  设 E 是维数为 n 的自由 A-模。~ \(x \in \bigwedge(\mathrm{E})\) 与 \(y \in \bigwedge(\mathrm{E})\) 关于一个基 \(\{e\}\) （属于 \(\bigwedge^{n}(\mathrm{E})\) ）的回归积（记作 \(x \vee y\) ）是元素 \(\phi(\phi(x) \wedge \phi(y))\) ，同构 \(\phi\) 是相对于 e 的。该乘积仅在相差一个可逆因子的意义下有定义，该因子依赖于所选的基 \(\{e\}\) 。证明：若 x 是一个 p-向量，y 是一个 q-向量，则 \(x \vee y = 0\) 当 \(p + q < n\) ，且若 \(p + q \geqslant n\) ， \(x \vee y\) 是一个 \((p + q - n)\) -向量，使得
\end{enumerate}

\[
y \vee x = (- 1) ^ {(n - p) (n - q)} x \vee y.
\]

回归积是结合的、对加法是分配的，并在 \(\bigwedge(\mathrm{E})\) 上定义了一个含单位元的、同构于外代数结构的代数结构。把 \(x \vee y\) 的分量表示为 \(x\) 和 \(y\) 的分量（对 \(\mathbf{E}\) 的一个给定基）的函数。

\begin{enumerate}
\def\labelenumi{\arabic{enumi}.}
\setcounter{enumi}{16}
\tightlist
\item
  设 K 是一个交换域，E 是 K 上的一个向量空间。
\end{enumerate}

\begin{enumerate}
\def\labelenumi{(\alph{enumi})}
\item
  设 \(z\) 是一个 \(p\) -向量 \(\neq 0\) （在 \(E\) 上），并设 \(V_{z}\) 为 \(E\) 的由向量 \(x\) （满足 \(z \wedge x = 0\) ）组成的向量子空间。则 \(\dim(V_{z}) \leqslant p\) ；若 \((x_{i})_{1 \leqslant i \leqslant q}\) 是 \(V_{z}\) 中向量的一个自由系，则存在一个 \((p - q)\) -向量 \(v\) ，使得 \(z = v \wedge x_{1} \wedge \dots \wedge x_{q}\) 。若 \(E\) 是有限维的，则 \(V_{z} = \phi(M_{z}^{0})\) ，其中 \(M_{z}^{0}\) 是 \(E^{*}\) 中的子空间 \(M_{z}\) （与 \(z\) 相关联）的正交补。要使 \(z\) 为纯 \(p\) -向量，必要且充分的是 \(\dim(V_{z}) = p\) ，此时 \(V_{z} = M_{z}\) 。
\item
  设 \(v\) 为纯 \(p\) -向量， \(w\) 为纯 \(q\) -向量，而 \(V\) 和 \(W\) 为 \(E\) 的分别与 \(v\) 和 \(w\) 相关联的子空间。要使 \(V \subset W\) ，必要且充分的是 \(q \geqslant p\) 且存在一个 \((q - p)\) -向量 \(u\) ，使得 \(w = u \wedge v\) 。要使 \(V \cap W = \{0\}\) ，必要且充分的是 \(v \wedge w \neq 0\) ；此时子空间 \(V + W\) 与纯 \((p + q)\) -向量 \(v \wedge w\) 相关联。
\item
  设 \(u\) 和 \(v\) 为两个纯 \(p\) -向量，且 \(U\) 和 \(V\) 为分别与 \(u\) 和 \(v\) 相关联的子空间。要使 \(u + v\) 为纯的，必要且充分的是 \(\dim(U \cap V) \geqslant p - 1\) 。
\end{enumerate}

\begin{enumerate}
\def\labelenumi{\arabic{enumi}.}
\setcounter{enumi}{17}
\item
  设 \(z = \sum_{H} \alpha_{H} e_{H}\) 为 n 维向量空间 E 的一个非零纯 p-向量，用它关于 E 的任意基 \((e_{i})_{1 \leqslant i \leqslant n}\) 的分量表示。设 G 为 \([1, n]\) 的具有 p 个元素的子集，使得 \(\alpha_{G} \neq 0\) ；令 \((i_{h})_{1 \leqslant h \leqslant p}\) 为 G 中指标按递增顺序排列的序列，而 \((j_{k})_{1 \leqslant k \leqslant n - p}\) 为 G 的补集 \(G'\) 中指标按递增顺序排列的序列。对每对有序指标 \((h, k)\) （满足 \(1 \leqslant h \leqslant p\) , \(1 \leqslant k \leqslant n - p\) ），令 \(\beta_{hk}\) 为 z 的分量 \(\alpha_{H}\) （对应于 \([1, n]\) 的子集 H），该子集 H 由 G 中异于 \(i_{h}\) 的 p - 1 个指标以及指标 \(j_{k} \in G'\) 组成；设 X 为具有 p 行和 n - p 列的矩阵 \((\beta_{hk})\) 。给定 \([1, n]\) 的任意 p 元子集 L，使得 \(L \cap G\) 具有 \(q \geqslant 1\) 个元素，证明 \((\alpha_{G})^{q-1} \alpha_{L}\) 除符号外等于 X 的 q 阶子式，该子式由满足 \(i_{h} \in G \cap G\) 的指标 h 的行和满足 \(j_{k} \in L \cap G\) 的指标 k 的列组成。（把 z 写成 \(\alpha_{G} y_{1} \wedge y_{2} \wedge \cdots \wedge y_{p}\) 的形式，其中向量 \(y_{i}\) 使得：在以 \(y_{i}\) 为列的 n 行 p 列矩阵 Y 中，由指标属于 G 的行组成的子矩阵是单位矩阵。）
\item
  在有限维向量空间 E 上，设 \(z = x_{1} \wedge x_{2} \wedge \cdots \wedge x_{p}\) 为一个纯p-向量 \(\neq 0\) ；令 \(v^{*}\) 为一个q-形式（ \(q \leqslant p\) ) 且 \(u^{*}\) 的任意一个元素 \(\bigwedge(\mathrm{E}^{*})\) 。证明
\end{enumerate}

\[
(u ^ {*} \wedge v ^ {*}) \lrcorner z = (- 1) ^ {q (q - 1) / 2} \sum_ {\mathrm{H}} \rho_ {\mathrm{H}, \mathrm{K}} \langle v ^ {*}, x _ {\mathrm{H}} \rangle (u ^ {*} \lrcorner x _ {\mathrm{K}})
\]

其中 H 遍历所有子集的集合 \([1, p]\) 具有q个元素，K是H在 \([1, p]\) 和 \(x_{H}\) （相应地 \(x_{K}\) )表示 \(x_{i}\) 的集合，使得 \(i \in H\) （相应地 \(i \in K\) ).

\begin{enumerate}
\def\labelenumi{\arabic{enumi}.}
\setcounter{enumi}{19}
\item
  设 E 为 n 维向量空间，z 为 E 上的纯 p-向量，V 为 E 的与 z 相关联的向量子空间， \(u^{*}\) 为 \(E^{*}\) 上的一个纯 q-向量， \(W'\) 为 \(E^{*}\) 的与 \(u^{*}\) 相关联的子空间，而 W 为 E 的、维数为 n - q 且正交于 \(W'\) 的子空间。若 q \textless{} p，则 \(u^{*} \lrcorner z\) 是 E 上的一个纯 (p - q)-向量，它当 \(\dim(V \cap W) > p - q\) 时为零；若 \(\dim(V \cap W) = p - q\) ，则 \(V \cap W\) 与 \(u^{*} \lrcorner z\) 相关联。
\item
  \begin{enumerate}
  \def\labelenumii{(\alph{enumii})}
  \tightlist
  \item
    直接证明（不使用同构 \(\phi_p\) ）：每个 \((n - 1)\) -向量（在 \(n\) 维向量空间上的）都是纯的。
  \end{enumerate}
\end{enumerate}

\begin{enumerate}
\def\labelenumi{(\alph{enumi})}
\setcounter{enumi}{1}
\tightlist
\item
  设 A 是域 K 上的一个交换代数，其基由单位元 1 和三个元素 \(a_{1}\) , \(a_{2}\) , \(a_{3}\) 组成，它们彼此之间的乘积为零。设 E 为 A-模 \(\mathbf{A}^3\) ，而 \((e_i)_{1\leqslant i\leqslant 3}\) 为其典范基；证明双向量
\end{enumerate}

\[
a _ {1} (e _ {2} \land e _ {3}) + a _ {2} (e _ {3} \land e _ {1}) + a _ {3} (e _ {1} \land e _ {2})
\]

不是纯的。

\begin{enumerate}
\def\labelenumi{\arabic{enumi}.}
\setcounter{enumi}{21}
\tightlist
\item
  设 E 为一个有序集，其中每个区间 \([a, b]\) 都是有限的。设 S 为 \(E \times E\) 的由有序对 \((x, y)\) （满足 \(x \leqslant y\) ）组成的子集。证明在 A-模 \(\mathbf{A}^{(\mathrm{s})}\) （由 S 中元素的形式线性组合构成）上，通过取余积定义了一个余结合的余代数结构
\end{enumerate}

\[
c ((x, y)) = \sum_ {x \leqslant z \leqslant y} (x, z) \otimes (z, y).
\]

对于这个余代数，线性形式 \(\gamma\) 使得

\[
\gamma ((x, y)) = 0 \quad \text { if } \quad x \neq y \text { in   E }, \quad \gamma ((x, x)) = 1
\]

是一个余单位。

\begin{enumerate}
\def\labelenumi{\arabic{enumi}.}
\setcounter{enumi}{22}
\tightlist
\item
  设 C 是环 A 上的一个余代数。C 的一个子 A-模 S 称为 C 的子余代数，如果 \(c(S) \subset \operatorname{Im}(S \otimes S)\) 。
\end{enumerate}

\begin{enumerate}
\def\labelenumi{(\alph{enumi})}
\item
  若 C 是余结合的（相应地，余交换的），则其每个子余代数也是。若 \(\gamma\) 是 C 的一个余单位，则 \(\gamma\) 在 C 的任意子余代数上的限制是一个余单位。
\item
  余代数 \(\mathbf{A}^{(\mathbf{X})}\) （第 1 号，例 4）的子余代数有哪些？
\item
  若 \((\mathbf{S}_{\alpha})\) 是 \(\mathbf{C}\) 的一个子余代数族，则 \(\sum_{\alpha} \mathbf{S}_{\alpha}\) 是 \(\mathbf{C}\) 的一个子余代数（包含所有 \(\mathbf{S}_{\alpha}\) 的最小者）。
\end{enumerate}

\begin{enumerate}
\def\labelenumi{\arabic{enumi}.}
\setcounter{enumi}{23}
\tightlist
\item
  设 C 是环 A 上的一个余代数。C 的一个子 A-模 S 称为右（相应地，左）余理想，如果 \(c(S) \subset \operatorname{Im}(S \otimes C)\) （相应地， \(c(S) \subset \operatorname{Im}(C \otimes S)\) ）。
\end{enumerate}

\begin{enumerate}
\def\labelenumi{(\alph{enumi})}
\item
  子余代数（习题 23）既是右余理想又是左余理想。反之，若 \(\mathbf{A}\) 是一个域，则结论成立。
\item
  每个右（相应地，左）余理想的并都是右（相应地，左）余理想。
\item
  如果 S 是 C 的右（相应地，左）余理想，则正交补 \(S^{0}\) 是对偶代数 \(C^{*}\) 的一个右（相应地，左）理想。
\item
  若 A 是一个域，且 \(\mathfrak{I}'\) 是对偶代数 \(\mathbf{C}^*\) 的一个右（相应地，左）理想，则正交 \(\mathfrak{I}''^0\) （在 \(\mathbf{C}\) 中）是 \(\mathbf{C}\) 的一个右（相应地，左）余理想。（用反证法论证：假设存在 \(x \in \mathfrak{I}''^0\) ，使得 \(c(x) \notin \mathfrak{I}''^0 \otimes \mathbf{C}\) ；写成 \(c(x) = \sum_{i} y_i \otimes z_i\) ，其中 \(z_i\) 线性无关，并证明这将导致存在 \(y' \in \mathfrak{I}'\) 和 \(z' \in \mathbf{C}^*\) ，使得 \(\langle y' z', x \rangle \neq 0\) 。）
\end{enumerate}

由此推出：对于 C 的一个向量子空间 S，要使 S 成为右（相应地，左）余理想，当且仅当 \(S^{0}\) 是 \(C^{*}\) 的一个右（相应地，左）理想。

\begin{enumerate}
\def\labelenumi{(\alph{enumi})}
\setcounter{enumi}{4}
\tightlist
\item
  由 (d) 推出：若 A 是一个域，则 C 的任意一族右余理想（相应地，左余理想，相应地，子余代数）的交仍是右余理想（相应地，左余理想，相应地，子余代数）。要使 S 成为 C 的子余代数，必要且充分的是其正交补 \(S^{0}\) 是代数 \(C^{*}\) 的一个双边理想。
\end{enumerate}

\begin{enumerate}
\def\labelenumi{\arabic{enumi}.}
\setcounter{enumi}{24}
\tightlist
\item
  设 C 是环 A 上的一个余代数。C 的一个子 A-模 S 称为余理想，如果 \(c(\mathbf{S}) \subset \operatorname{Im}(\mathbf{S} \otimes \mathbf{C}) + \operatorname{Im}(\mathbf{C} \otimes \mathbf{S})\) 。若 C 是有余单位的，其余单位为 \(\gamma\) ，则余理想 S 称为余零的，如果 \(\gamma(x) = 0\) （对 \(x \in S\) ）。
\end{enumerate}

\begin{enumerate}
\def\labelenumi{(\alph{enumi})}
\item
  右或左余理想都是余理想。余理想的任意和仍是余理想。
\item
  若 S 是 C 的余理想，则其正交补 \(S^{0}\) 是对偶代数 \(C^{*}\) 的一个子代数。若 C 是有余单位的且 S 是余零的，则 \(S^{0}\) 是含单位元的。
\item
  设 A 是一个域，C 是一个有余单位的余代数。证明若 \(E'\) 是 \(C^{*}\) 的一个含单位元的子代数（与 \(C^{*}\) 有相同的单位元），则 C 中的正交 \(E'^{0}\) 是一个余零余理想。
\item
  如果 S 是 C 的一个余理想，证明在 C/S 上存在且仅存在一个余代数结构，使得典范映射 \(p: C \rightarrow C/S\) 是一个余代数态射。如果 C 是有余单位的且 S 是余零的，则 C/S 是有余单位的。
\item
  对于每个余代数态射 \(f\colon\mathbf{C}\to\mathbf{C}^{\prime},\mathbf{S}=\operatorname{Ker}(f)\) 是 \(\mathbf{C}\) 的余理想， \(\mathbf{S}^{\prime}=\operatorname{Im}(f)\) 是 \(\mathbf{C}^{\prime}\) 的一个子余代数，并且在典范分解
\end{enumerate}

\[
\mathrm{C} \xrightarrow {p} \mathrm{C} / \mathrm{S} \xrightarrow {g} \mathrm{S} ^ {\prime} \xrightarrow {j} \mathrm{C} ^ {\prime}
\]

中 \(f, g\) 是一个余代数同构。

\begin{enumerate}
\def\labelenumi{(\arabic{enumi})}
\setcounter{enumi}{25}
\tightlist
\item
  设 C 是交换域 K 上的一个余结合、有余单位的余代数。
\end{enumerate}

\begin{enumerate}
\def\labelenumi{(\alph{enumi})}
\item
  对于 C 的每一个向量子空间 E（它在合成律 \((u, x) \mapsto u \perp x\) 下是一个左 C*-模），这个 C*-模的左零化子是一个双边理想 \(S'\) （在代数 \(C^*\) 中）。证明若 E 在 K 上是有限维的，则 \(C^*/S'\) 是有限维的。
\item
  从 (a) 推出，对每个元素 \(x \in \mathbb{C}\) ， \(\mathbb{C}\) 的由 \(x\) 生成的子余代数是有限维的。（注意， \(\mathbb{C}^*\) -模 \(\mathbb{E}\) （由 \(x\) 生成的）是有限维的，并且若 \(\mathfrak{S}'\) 是其左零化子，则 \(x \in \mathfrak{S}'^0\) ，即 \(\mathfrak{S}'\) 在 \(\mathbb{C}\) 中的正交。）
\end{enumerate}

\begin{enumerate}
\def\labelenumi{(\arabic{enumi})}
\setcounter{enumi}{26}
\tightlist
\item
  设 B 是交换域 K 上的一个结合含幺代数，并设 \(m: B \otimes_{K} B \to B\) 是定义 B 上乘法的 K-线性映射。张量积 \(B^{*} \otimes_{K} B^{*}\) （其中 \(B^{*}\) 是向量空间 B 的对偶向量空间）典范地等同于 \((\mathbf{B} \otimes_{\mathbf{K}} \mathbf{B})^{*}\) 的一个向量子空间（II，§7，第 7 号，命题 16）。
\end{enumerate}

\begin{enumerate}
\def\labelenumi{(\alph{enumi})}
\item
  证明对于元素 \(w \in \mathbf{B}^*\) ，以下条件等价：（1） \({}^t m(w) \in \mathbf{B}^* \otimes \mathbf{B}^*\) ；(2) 元素 \(x \lrcorner w\) （其中 \(x\) 遍历 \(\mathbf{B}\) ）组成的集合是 \(\mathbf{B}^*\) 的一个有限维子空间；(3) 元素 \(w \lfloor x\) （其中 \(x\) 遍历 \(\mathbf{B}\) ）组成的集合是 \(\mathbf{B}^*\) 的一个有限维子空间；(4) 元素 \(x \lrcorner w \lfloor y\) （其中 \(x\) 和 \(y\) 遍历 \(\mathbf{B}\) ）组成的集合包含在 \(\mathbf{B}^*\) 的一个有限维子空间中。
\item
  设 \(\mathbf{B}' \subset \mathbf{B}^*\) 为 \(w \in \mathbf{B}^*\) 中满足
\end{enumerate}

(a) 中的等价条件者的集合。证明 \({}^t m(\mathbf{B}') \subset \mathbf{B}' \otimes_{\mathbb{K}} \mathbf{B}'\) 。（对于 \(w \in \mathbf{B}'\) ，记 \({}^t m(w) = \sum_i u_i \otimes v_i\) ，其中例如 \(u_i\) 在 \(\mathbf{B}^*\) 中线性无关；然后证明 \(v_i \in \mathbf{B}'\) （对所有 \(i\) ），利用 \(\mathbf{B}\) 的结合性并通过反证法论证。）\(\mathbf{B}'\) 因此是包含在 \(\mathbf{B}^*\) 中的最大余代数，对其而言 \({}^t m\) 是余积。\(\mathbf{B}'\) 称为代数 \(\mathbf{B}\) 的对偶余代数。当 \(\mathbf{B}\) 是有限维的时， \(\mathbf{B}' = \mathbf{B}^*\) 。

\begin{enumerate}
\def\labelenumi{(\alph{enumi})}
\setcounter{enumi}{2}
\item
  证明：如果 \(\mathbf{B}\) 是 \(\mathbf{K}\) 上无限维的交换域，则 \(\mathbf{B}' = \{0\}\) 。（注意，元素 \(x \perp w\) （对固定的 \(w \in \mathbf{B}^*\) ，\(x\) 遍历 \(\mathbf{B}\) ）组成的集合的正交补是 \(\mathbf{B}\) 的一个理想。）
\item
  设 \(\mathbf{B}_1, \mathbf{B}_2\) 是两个 K-代数，且 \(f: \mathbf{B}_1 \to \mathbf{B}_2\) 是代数的一个 K-同态。证明 \(^t f(\mathbf{B}_2') \subset \mathbf{B}_1'\) ，并且 \(^t f\) 限制在 \(\mathbf{B}_2'\) 上是从 \(\mathbf{B}_2'\) 到 \(\mathbf{B}_1'\) 中的余代数同态。
\item
  若 \(\mathbf{C}\) 是 \(\mathbf{K}\) 上的一个余结合、有余单位的余代数，则典范单射 \(\mathbf{C} \to \mathbf{C}^{**}\) （向量空间 \(\mathbf{C}\) 到其二对偶中的）把 \(\mathbf{C}\) 映到 \((\mathbf{C}^*)'\) ——代数 \(\mathbf{C}^*\) （对偶于 \(\mathbf{C}\) ）的对偶余代数——的一个子空间上，并且是从 \(\mathbf{C}\) 到 \((\mathbf{C}^*)'\) 中的余代数同态。
\end{enumerate}

\BourbakiAppendixExercises

\begin{enumerate}
\def\labelenumi{\arabic{enumi}.}
\item
  证明：在交错 Cayley A-代数中（§ 2 第 4 号的记号）， \(\mathrm{T}((xy)z) = \mathrm{T}(x(yz))\) 。（把它归结为证明 \(a(x,y,z) = a(\bar{z},\bar{y},\bar{x})\) 。）
\item
  设 F 是一个交错 A-代数，其中关系 2x = 0 蕴含 x = 0。证明：对 F 中所有 x, y, z， \(x(yz)x = (xy)(zx)\) （Moufang 恒等式）。（在该恒等式中
\end{enumerate}

\[
(x y) z + y (z x) = x (y z) + (y z) x
\]

依次把 \(y\) 换成 \(xy\) ，把 \(z\) 换成 \(zx\) 。

\begin{enumerate}
\def\labelenumi{\arabic{enumi}.}
\setcounter{enumi}{2}
\tightlist
\item
  设 \(\mathbf{F}\) 是一个交错 Cayley A-代数，其中关系 \(2x = 0\) 蕴含 \(x = 0\) 。证明如果 \(x \in \mathbf{F}\) 是可逆的，则对所有 \(y, z, t\) （在 \(\mathbf{F}\) 中），
\end{enumerate}

\[
(\mathrm{N} (x) + \mathrm{N} (y)) (\mathrm{N} (z) + \mathrm{N} (t)) = \mathrm{N} (x \bar {z} + y \bar {t}) + \mathrm{N} (x t - (x z) (x ^ {- 1} y))
\]

（利用习题 1 和习题 2）。
